\documentclass[11pt,a4paper]{article}

% ========================
% ESSENTIAL PACKAGES
% ========================

% Language and encoding
\usepackage[utf8]{inputenc}
\usepackage[T1]{fontenc}
\usepackage[english]{babel}
%\usepackage{natbib}

% Math packages
\usepackage{amsmath}        % Enhanced math environments
\usepackage{amssymb}        % Additional math symbols
\usepackage{amsthm}         % Theorem environments
\usepackage{amsfonts}       % Additional math fonts
\usepackage{mathtools}      % Extension to amsmath
\usepackage{bm}             % Bold math symbols
\usepackage{dsfont}         % Double-struck fonts (for sets like ℝ, ℕ)
\usepackage{mathrsfs}       % Script fonts for math

% Graphics and figures
\usepackage{graphicx}       % Include graphics
\usepackage{float}          % Better float placement
\usepackage{subcaption}     % Subfigures
\usepackage{tikz}           % Drawing diagrams
\usepackage{pgfplots}       % Plotting
\pgfplotsset{compat=1.17}

% Page layout and formatting
\usepackage[margin=1in]{geometry}
\usepackage{fancyhdr}       % Custom headers and footers
\usepackage{titlesec}       % Customize section titles
\usepackage{enumitem}       % Customize lists
\usepackage{parskip}        % Paragraph spacing

% Colors and highlighting
\usepackage{xcolor}
\usepackage{mdframed}       % Framed environments
\usepackage{tcolorbox}      % Colored boxes

% Tables
\usepackage{array}          % Extended array environments
\usepackage{booktabs}       % Professional tables
\usepackage{multirow}       % Multi-row cells

% References and links
\usepackage{hyperref}       % Hyperlinks
\usepackage{cleveref}       % Smart cross-references

% Additional useful packages
\usepackage{cancel}         % Cancel terms in equations
\usepackage{siunitx}        % SI units
\usepackage{listings}       % Code listings (if needed)

\usepackage{authblk}

\usepackage[authoryear]{natbib}

% ========================
% CUSTOM COMMANDS
% ========================

% Common math sets
\newcommand{\R}{\mathbb{R}}
\newcommand{\N}{\mathbb{N}}
\newcommand{\Z}{\mathbb{Z}}
\newcommand{\Q}{\mathbb{Q}}
\newcommand{\C}{\mathbb{C}}

\newcommand{\ks}{\mathrm{KS}}

% Common operators
\DeclareMathOperator{\lcm}{lcm}
\DeclareMathOperator{\Tr}{Tr}
\DeclareMathOperator{\rank}{rank}
%\DeclareMathOperator{\span}{span}
%\DeclareMathOperator{\null}{null}

% Shortcuts for common expressions
\newcommand{\abs}[1]{\left|#1\right|}
\newcommand{\norm}[1]{\left\|#1\right\|}
\newcommand{\inner}[2]{\left\langle #1, #2 \right\rangle}
\newcommand{\ceil}[1]{\left\lceil #1 \right\rceil}
\newcommand{\floor}[1]{\left\lfloor #1 \right\rfloor}

% Derivatives and integrals
\newcommand{\dd}[1]{\,\mathrm{d}#1}
\newcommand{\dv}[2]{\frac{\mathrm{d}#1}{\mathrm{d}#2}}
\newcommand{\pdv}[2]{\frac{\partial #1}{\partial #2}}

\renewcommand{\hat}{\widehat}
\renewcommand{\tilde}{\widetilde}
\renewcommand{\bar}{\overline}

\newcommand{\vct}{\boldsymbol}
\newcommand{\tr}{\mathrm{tr}}

% ========================
% THEOREM ENVIRONMENTS
% ========================

\theoremstyle{definition}
\newtheorem{definition}{Definition}[section]
\newtheorem{example}{Example}[section]
\newtheorem{exercise}{Exercise}[section]

\theoremstyle{plain}
\newtheorem{theorem}{Theorem}[section]
\newtheorem{lemma}{Lemma}[section]
\newtheorem{corollary}{Corollary}[section]
\newtheorem{proposition}{Proposition}[section]
\newtheorem{assumption}{Assumptio}[section]

\theoremstyle{remark}
\newtheorem{remark}{Remark}[section]
\newtheorem{note}{Note}[section]

\newcommand{\supp}{\mathrm{supp}}
\newcommand{\op}{\mathrm{op}}

\newcommand{\mX}{\mathcal X}
\newcommand{\mD}{\mathcal D}
\newcommand{\bP}{\mathbb P}
\newcommand{\lp}{L^2(\bP)}
\newcommand{\lx}{L^2(\mX)}
\newcommand{\ltwo}{L^2(\mX)}
\newcommand{\ud}{\mathrm d}
\newcommand{\mH}{\mathcal H}
\newcommand{\mA}{\mathcal A}
\newcommand{\mB}{\mathcal B}

\newcommand{\diag}{\mathrm{diag}}

\newcommand{\mI}{\mathcal I}

\newcommand{\iso}{\mathsf{iso}}
\newcommand{\ani}{\mathsf{ani}}
\newcommand{\anitr}{\mathsf{anitr}}
\newcommand{\adptr}{\mathsf{adptr}}

\newcommand{\und}{\underline}
\newcommand{\sgn}{\mathrm{sgn}}
\newcommand{\pd}[1]{
	{ \color{blue} PD: #1}
}

% ========================
% COLORED BOXES
% ========================

% Important theorem box
\newtcolorbox{importantbox}{
	colback=blue!5!white,
	colframe=blue!75!black,
	title=Important
}

% Warning/Note box
\newtcolorbox{notebox}{
	colback=yellow!10!white,
	colframe=orange!75!black,
	title=Note
}

% Example box
\newtcolorbox{examplebox}{
	colback=green!5!white,
	colframe=green!75!black,
	title=Example
}

% ========================
% DOCUMENT SETUP
% ========================

% Page style
\pagestyle{fancy}
\fancyhf{}
\rhead{\thepage}
\lhead{\leftmark}


% Hyperlink setup
\hypersetup{
	colorlinks=true,
	linkcolor=blue,
	filecolor=magenta,      
	urlcolor=cyan,
	citecolor=red
}

\title{Honest and Adaptive Pointwise Confidence Intervals in Linear Functional Regression: An Anisotropic Kernel Ridge Regression Approach}
\author[1]{Paromita Dubey \footnote{Author names listed in alphabetical order.}}
\author[2]{Zhiyuan Tang}
\author[2]{Yining Wang}
\affil[1]{Marshall Business School, University of Southern California}
\affil[2]{Naveen Jindal School of Management, University of Texas at Dallas}


\begin{document}

\maketitle

\begin{abstract}
Abstract here.
\end{abstract}

\section{Introduction}

We consider the classical model of linear functional regression with data $\{X_i,Y_i\}_{i=1}^n$, where $X_1,\cdots,X_n\sim P_X$ are independently and identically distributed \emph{functions} 
whose domain is a non-empty compact interval $I\subseteq\mathbb R$, and $\{Y_i\}_{i=1}^n$ are scalar response variables, modeled by
\begin{equation}
\mathbb E[Y_i|X_i] = \eta_0(X_i) := \gamma_0 + \int_I X_i(t)\beta_0(t)\ud t,
\label{eq:model}
\end{equation}
where $\gamma_0\in\mathbb R$ is the intercept term and $\beta_0: I\to\mathbb R$ is the linear functional, both being unknown parameters.
Many well-known methods have been proposed to estimate $\gamma_0$ and $\beta_0$ from data, such as functional principal component analysis (fPCA) approaches \citep[Chapter 8]{ramsay2005functional},
\citep{cai2006prediction} and methods based on reproducing kernel Hilbert spaces \citep{yuan2010reproducing,cai2012minimax}.

In this paper, we consider the question of constructing \emph{pointwise} confidence intervals on $\eta_0(x)$ which are valuable in quantifying the uncertainty in estimates of the functional model
when it is applied to a future data curve $x$. We study the construction of \emph{honest} confidence intervals, which are functions $\varphi$ mapping a data curve $x:I\to\mathbb R$ to a compact interval $\varphi(x)\subseteq\mathbb R$
such that
\begin{equation}
\liminf_{n\to\infty}\inf_{\eta_0\in\Theta} \Pr\left[\eta_0(x)\in\varphi(x)\right] \geq 1-\alpha,\;\;\;\;\;\;\forall x\in\mathrm{supp}(P_X)
\label{eq:honest-ci}
\end{equation}
where $\alpha\in(0,1)$ is a pre-determined significance level,
$\supp(P_X)$ is the support of $P_X$, $\Theta$ is a fixed and known function class and the randomness in Eq.~(\ref{eq:honest-ci}) is over $\{(X_i,Y_i)\}_{i=1}^n$.
To understand how \emph{adative} an honest confidence interval is, we further study the worst-case mean-squared widths of constructed intervals
\begin{equation}
\sup_{\eta_0\in\Theta}\mathbb E_{\eta_0} \left[\int_{\supp(P_X)}\big|\varphi(x)\big|^2\ud P_X(x)\right],
\label{eq:adaptive}
\end{equation}
where $|\varphi(x)|$ is the Lesbesgue measure of the output interval $\varphi(x)$. We are particularly interested in how the above worst-case risk converges to zero as $n\to\infty$,
and how it compares with the mean-squared prediction errors
$$
\sup_{\eta_0\in\Theta}\mathbb E_{\eta_0} \left[\int_{\supp(P_X)}\int_I [x(t)(\hat\eta_n(t)-\eta_0(t)]^2\ud t\ud P_X(x)\right]
$$
for classical fPCA or RKHS estimators $\hat\eta$ of $\eta_0$.

\subsection{Our contributions}

\subsection{Related works}

\subsection{Notations}

For two sequences of real numbers $\{a_n\}$ and $\{b_n\}$, we write $a_n\lesssim b_n$, or equivalently $a_n=O(b_n)$, if there exists a numerical constant $C$ such that $|a_n|\leq C|b_n|$ for all $n$.
We write $a_n\gtrsim b_n$ or $a_n=\Omega(b_n)$ if $b_n\lesssim a_n$, and $a_n\asymp b_n$ or $a_n=\Theta(b_n)$ if both $a_n\lesssim b_n$ and $a_n\gtrsim b_n$ hold.
We denote $L_2(I):=\{x: I\to\mathbb R: \int_I x(t)^2\ud t<+\infty\}$ as the class of all square integrable functions on $I$.
For two functions $x,z\in L_2(I)$, $\langle x,z\rangle$ denotes the standard inner product of $L_2(I)$; that is, $\langle x,z\rangle = \int_I x(t)z(t)\ud t$.
For two real numbers $a,b\in\mathbb R$, we write $a\vee b=\max\{a,b\}$ and $a\wedge b=\min\{a,b\}$.

\section{Functional regression via kernel ridge regression}

In this section, we give an overview of an reproducing kernel Hilbert space (RKHS) approach towards linear functional regression, following the seminal works of \cite{cai2012minimax} which also lays the foundation of 
methodological and theoretical contributions of this paper.
To clarify, all results in this section are relatively well-known in the literature, and we are summarizing them for the purpose of completeness only.

Let $K:I\times I\to\mathbb R$ be a real, symmetric, square integrable, and non-negative definite function, known as the \emph{kernel function}.
There then exists a reproducing kernel Hilbert space $\mH\subseteq L_2(I)$ associated with inner product $\langle\cdot,\cdot\rangle_{\mH}$
such that the following reproducing property holds for any $f\in\mH$:
$$
f(t) = \langle K(t,\cdot), f\rangle_{\mH}\;\;\;\;\;\;\forall t\in I.
$$
Assuming $\beta_0\in\mH$. Then a standard kernel ridge regression (KRR) approach to estimate $\alpha_0,\beta_0$ from $\{X_i,Y_i\}_{i=1}^n$ is to solve the following penalized optimization problem
\begin{equation}
\hat\gamma_n,\hat\beta_n \in \arg\min_{\gamma\in\mathbb R,\beta\in\mH} \left\{\frac{1}{n}\sum_{i=1}^n (Y_i-\gamma-\langle X_i,\beta\rangle)^2 + \lambda_n \|\beta\|_\mH^2\right\},
\label{eq:krr}
\end{equation}
where $\lambda_n>0$ is a penalization parameter.
Let $\bar X=\frac{1}{n}\sum_{i=1}^n X_i$ and $\bar Y=\frac{1}{n}\sum_{i=1}^n Y_i$ be the sample averages. It is easy to verify that $\hat\gamma_n=\bar Y$ and 
\begin{equation}
\hat\beta_n\in\arg\min_{\beta\in\mH}\left\{\frac{1}{n}\sum_{i=1}^n (Y_i'-\langle X_i',\beta\rangle)^2 + \lambda_n\|\beta\|_\mH^2\right\}
\label{eq:krr-nointercept}
\end{equation}
where $Y_i'=Y_i-\bar Y$ and $X_i'=X_i-\bar X$.

\subsection{Reduction to an optimization problem in $L_2(I)$}

The work of \cite{cai2012minimax} makes the observation that Eq.~(\ref{eq:krr-nointercept}) can be equivalently reformulated into an infinite-dimensional optimization problem
involving the standard inner product and norm of $L_2(I)$ only, which is helpful for implementation and theoretical analysis purposes.
We summarize how such equivalence is built below.

By Mercer theorem, there exists an orthonormal basis $\{\phi_j\}_{j=1}^{\infty}$ of $L_2(I)$
and a sequence of non-increasing, strictly positive real numbers $\rho_1\geq\rho_2\geq\cdots$ such that $\phi_j\in\mH$ for all $j$, and furthermore
$$
\langle\phi_j,\phi_k\rangle = \int_I \phi_j(t)\phi_k(t)\ud t = \delta_{jk},\;\;\;\;\;\langle\phi_j,\phi_k\rangle_{\mH}= \frac{\delta_{jk}}{\rho_j},\;\;\;\;\;\;\int_I K(t,\cdot)\phi_j(t)\ud t=\rho_j\phi_j.
$$
where $\delta_{jk}=\vct 1\{j=k\}$ is the Kronecker delta notation.
Note that because $\mH$ is known, $\{\phi_j,\rho_j\}_{j=1}^{\infty}$ are also known.
Following notations in \citep{cai2012minimax}, we define linear operator $L_{K^{1/2}}:L^2(I)\to L^2(I)$ and its inverse operator $L_{K^{-1/2}}$ as follows:
$$
L_{K^{1/2}}(\phi_j) = \sqrt{\rho_j}\phi_j,\;\;\;\;\;\; L_{K^{-1/2}}(\phi_j) = \frac{1}{\sqrt{\rho_j}}\phi_j,\;\;\;\;\;\forall j\geq 1.
$$
In other words, for any $f=\sum_{j=1}^{\infty}\alpha_j\phi_j\in L^2(I)$ such that $\sum_j\alpha_j^2<+\infty$, $L_{K^{1/2}}(f)=\sum_{j=1}^{\infty}\alpha_j\sqrt{\rho_j}\phi_j$
and therefore $L_{K^{1/2}}f\in\mH$ because $\sum_{j=1}^{\infty}\langle L_{K^{1/2}}f,\phi_j\rangle^2/\rho_j <+\infty$.
In fact, \citep[Proposition 2]{yuan2010reproducing} shows that for $f\neq 0$ and $L_{K^{1/2}}f\neq 0$, $f\in L^2(I)$ if and only if $L_{K^{1/2}}f\in\mH$ and furthermore $\int_I f(t)^2\ud t=\|f\|_2^2=\|L_{K^{1/2}}f\|_\mH^2$.
The optimal solution $\hat\beta_n$ in Eq.~(\ref{eq:krr-nointercept}) can then be expressed as $\hat\beta_n = L_{K^{1/2}}\hat\theta_n$, where
\begin{equation}
\hat\theta_n \in \arg\min_{\theta\in L_2(I)}\left\{\frac{1}{n}\sum_{i=1}^n (Y_i'-\langle L_{K^{1/2}}X_i', \theta\rangle^2 + \lambda_n\|\theta\|_2^2\right\},
\label{eq:krr-nointercept-l2}
\end{equation}
where $\|\theta\|_2^2 =\int_I \theta(t)^2\ud t$ is the standard squared $L_2$ norm on $I$.
For brevity, we will also use the notation of $Z_i' = L_{K^{1/2}}X_i'$ throughout this paper.

\subsection{Kernel ridge regression via spectral decomposition}

The infinite-dimensional optimization problem in Eq.~(\ref{eq:krr-nointercept-l2}) can be reduced to a finite-dimensional problem and efficiently solved
using Gram matrices, known as the kernel trick. In this section, we show an alternative characterization of $\hat\theta_n$, representing it as a smoothed version of functional PCA
which is useful motivating our anisotropic estimators later.

Let
$$
T_n := \frac{1}{n}\sum_{i=1}^n Z_i'\otimes Z_i'
$$
be the sample covariance operator, where $Z_i'=L_{K^{1/2}}X_i'$. By the spectral theorem, there exists orthonormal $\{\hat\psi_j\}_{j=1}^n$ in $L_2(I)$ and $\hat s_1\geq \hat s_2\geq\cdots\geq \hat s_n\geq 0$ such that
$$
T_n = \sum_{j=1}^n \hat s_j\hat \psi_j\otimes\hat \psi_j.
$$
The estimator $\hat\theta_n$ can then be expressed as
\begin{equation}
\hat\theta_n = (T_n+\lambda_n I)^{-1}\frac{1}{n}\sum_{i=1}^n Y_i'Z_i' = \frac{1}{n}\sum_{i=1}^n Y_i'\sum_{j=1}^{\infty}\frac{\hat\nu_j(Z_i')}{\hat s_j+\lambda}\hat\psi_j,
\label{eq:krr-nointercept-fpca}
\end{equation}
where $\hat\nu_j(Z_i') = \langle Z_i', \hat\psi_j\rangle=\int_I Z_i'(t)\hat\psi_j(t)\ud t$.

\section{Assumptions}

In this section we list and discuss assumptions that we impose throughout this paper. For some results, additional assumptions are needed: 
we will remark on them when such additional assumptions are introduced specifically for selected theorems or lemmas.

We first introduce mild regularity conditions to make sure the problem is well-posed.
\begin{enumerate}
\item[(A1)] Noise variables $\varepsilon_i:=Y_i-\eta_0(X_i)$ are independently and identically distributed, centered random variables with $\mathbb E[\varepsilon_i^2]=\sigma_\varepsilon^2<+\infty$
and $\mathbb E[|\varepsilon_i|^3]=M_\varepsilon < +\infty$; \pd{We need to add the estimation of $\sigma_\varepsilon^2$, as it enters directly into the construction of the CIs.}
\item[(A2)] $\beta_0\in\mH$ and $\|\beta_0\|_\mH\leq 1$, which is equivalent to $\|\theta_0\|_2=\|L_{K^{-1/2}}\beta_0\|_2 =(\int_I \theta_0(t)^2\ud t)^{1/2} \leq 1$.
\item[(A3)] For $x\sim P_X$, $\|L_{K^{1/2}}x\|_2= (\int_I (L_{K^{1/2}}x)(t))^2\ud t)^{1/2}\leq 1$ almost surely. We also assume that $x\neq 0$ and $L_{K^{1/2}}x\neq 0$ almost surely, to avoid degenerate situations.
\end{enumerate}

We remark that in Assumptions (A2) and (A3) we use one as the upper bounds to $\|\beta_0\|_\mH$ and $\|x\|_2$ for convenience only. Our constructed honest confidence bands do not require such knowledge
and they remain valid with 1 being replaced by any other numerical constants.

We use $\Theta_2(\mH)$ to denote the class of all regression functionals $\beta_0$ that satisfy (A2). Note that $\Theta_2(\mH)$ depends on the selection of the RKHS $\mH$, but not on the covariate distribution $P_X$.

Let $Z:=L_{K^{1/2}}X$ be a random function and $P_Z$ be the distribution of $Z$, which is induced by both $P_X$ and $K$. Consider the population covariance operator
$$
T := \mathbb E_{P_Z}\left[Z\otimes Z\right].
$$
By the spectral theorem, there exists an orthonormal basis $\{\psi_j\}_{j=1}^{\infty}$ of $L_2(I)$ and $s_1\geq s_2\geq\cdots\geq 0$ such that
\begin{equation}
T = \sum_{j=1}^{\infty} s_j\psi_j\otimes\psi_j.
\label{eq:specT}
\end{equation}
Note that, in general, because the population covariance operator $T$ is unknown, $\{\psi_j,s_j\}$ are also unknown.

The following assumptions on imposed directly on the structure of $\{\psi_j\}_j$ and $\{s_j\}_j$.
\begin{enumerate}
\item[(B1)] There exists numerical constants $r>1/2$ and $0<c_s\leq C_s<+\infty$ such that $c_s j^{-2r}\leq s_j\leq C_sj^{-2r}$ for all $j\geq 1$;
this also implies that $T$ is invertible;
\item[(B2)] There exists a numerical constant $C_p<+\infty$ such that for every $f\in L_2(I)$, $\mathbb E[\langle Z,f\rangle^4] \leq C_p(\mathbb E[\langle Z,f\rangle^2])^2$;
\item[(B3)] There exists numerical constants $\kappa\in[0,1/2)$ and $C_\nu<+\infty$ such that for every $j\geq 1$, $|\langle Z,\psi_j\rangle| \leq C_\nu n^{\kappa}\cdot \sqrt{s_j}$ almost surely.
\end{enumerate}

Assumption (B2) is also imposed, in the form of Assumption (b) in \citep{yuan2010reproducing}. It is satisfied when $X$ is a Gaussian process.

We remark that, in Assumption (B3), the smaller $\kappa$ the stronger the assumption is. For the strongest version of $\kappa=0$, it essentially assumes $|\langle Z,\psi_j\rangle|\lesssim j^{-r}$,
which coincides with coefficient decay assumptions adopted in many fPCA analysis.
{\color{red}In particular, the results in \cite{cai2006prediction} shows that in this case ($\kappa=0$ in (B3)) fPCA achieves an MSE rate of $n^{-(2r-1)/(2r+1)}$ assuming no decay of $|\langle\theta_0,\psi_j\rangle|$.}

%
%To analyze truncated estimators, we also introduce the following assumption which is strictly weaker than Assumption (B3):
%\begin{enumerate}
%\item[(B3')] For some $q\in(0,1/2)$, there exists numerical constants $\kappa<1/2$ and $C_\nu<+\infty$ such that for any $\delta>0$, there exists $n_0\in\mathbb N$ depending on $\delta$ such that
%for every $n\geq n_0$ and $1\leq j\leq n^{q+\delta}$, $|\langle Z,\psi_j\rangle|\leq C_\nu n^{\kappa}\cdot\sqrt{s_j}$.
%\end{enumerate}
%
%Given Assumption (B1), the following condition is a sufficient condition that implies Assumption (B3'): for some $a<\kappa/q$, it holds almost surely that
%$$
%\big|\langle Z,\psi_j\rangle\big|\lesssim j^{-r+a}\;\;\;\;\;\;\forall j\geq 1.
%$$

{\color{red}In light of Assumption (B3) which is weaker for larger $n$, we allow the distribution $P_X$ (and therefore the population covariance operator $T$) to vary for different $n$,
so that more distributions $P_X$ satisfy Assumption (B3) for larger values of $n$.
When necessary, we use the notation $P_X^n$ to emphasize on this fact.
We also use $P_Z^n$ to denote the distribution of $Z=L_{K^{1/2}}X$ when $X\sim P_X$.}

{\color{red}It should also be remarked that the assumption (B3) with $\kappa=0$ would greatly simplify the problem, because isotropic weight choices are automatically adaptive.
Similar phenomenon revealed in \citep{liu2025scalable}.}

\section{Under-smoothed isotropic kernel regression}\label{sec:isotropic}

As a warm-up exercise and a backup method with minimal assumptions, we first consider honest confidence intervals constructed by \emph{under-smoothing} an isotropic kernel ridge estimator.
Let $\{\lambda_n\}_{n=1}^{\infty}$ be a sequence of positive penalty parameters that satisfies
\begin{equation}
\lambda_n n \to 0,\;\; \lambda_n n^{1+\delta}\to \infty\;\;\;\;\;\;\text{as}\;\;n\to\infty
\label{eq:lambdan-1}
\end{equation}
for any $\delta>0$.
Let $\hat\gamma_n = \bar Y$ and $\hat\beta_n^\iso=L_{K^{1/2}}\hat\theta_n^\iso$ where 
$$
\hat\theta_n^\iso = (T_n+\lambda_n I)^{-1}\frac{1}{n}\sum_{i=1}^n Y_i'Z_i',
$$
with $Y_i'=Y_i-\bar Y$, $Z_i'=L_{K^{1/2}}(X_i-\bar X)$ and $T_n=\frac{1}{n}\sum_{i=1}^n Z_i'\otimes Z_i'$.
This is an \emph{under-smoothed} kernel ridge regression estimator with $\lambda_n=o(n^{-1})$, while the optimal penalty parameter choice for estimating $\eta_0$ in mean-squared errors 
is $\lambda_n\asymp n^{-2r/(2r+1)}$, per results in \cite{cai2012minimax}.

For a test function $x$ in the support $P_X^n$, an honest and adaptive confidence interval of $\eta(x)$ at significance level $\alpha\in(0,1)$ is constructed as
$$
\varphi^\iso(x) = [\hat\eta_n^\iso(x)\;\pm\;\tau_{\alpha/2}\hat\sigma_n^\iso(L_{K^{1/2}}x)]
$$
where $\tau_{\alpha/2}$ is the $(1-\alpha/2)$-quantile of the standard Normal distribution, and
$$
\hat\eta_n^\iso(x) = \hat\gamma_n + \int_I x(t)\hat\beta_n^\iso(t)\ud t,\;\;\;\;\;\;\hat\sigma_n(z)^\iso = \frac{\sigma_\varepsilon}{\sqrt{n}}\sqrt{\langle z, (T_n+\lambda_n I)^{-1}z\rangle}.
$$
%To facilitate our analysis, we also introduce the following quantity $\tilde z_n(z)$ defined as
%$$
%\tilde\sigma_n(z) := \frac{\sigma_\varepsilon}{\sqrt{n}}\sqrt{\langle z, (T_n+\lambda_n I)^{-1}T_n(T_n+\lambda_n I)^{-1}z\rangle},
%$$
%which is smaller than $\hat\sigma_n(z)$ almost surely.

\subsection{Honesty of confidence intervals}

For the purpose of our analysis, 
we also introduce the following quantity $\tilde z_n(z)$ defined as
$$
\tilde\sigma_n(z) := \frac{\sigma_\varepsilon}{\sqrt{n}}\sqrt{\langle z, (T_n+\lambda_n I)^{-1}T_n(T_n+\lambda_n I)^{-1}z\rangle},
$$
which is smaller than $\hat\sigma_n^\iso(z)$ almost surely.

Let $\mX := \cup_{n=1}^{\infty}\cap_{m=n}^{\infty}\supp(P_X^m)$ denote the class of all functions which belong to the support of $P_X^n$ for all sufficiently large $n$.

\begin{theorem}
Suppose Assumptions (A1), (A2), (A3), (B1), (B2), (B3) hold, where (B3) holds with $\kappa<\frac{1}{6}-\frac{1}{4r}$. Suppose also that Eq.~(\ref{eq:lambdan-1}) holds. Then 
for any $x\in\supp(P_X)$, one can decompose the pointwise prediction error at $x$ as $\hat\eta_n^\iso(x)-\eta_0(x)=b_n(x)+v_n(x)$ such that, as $n\to\infty$, 
\begin{align*}
\frac{b_n(x)}{\hat\sigma_n^\iso(z)}\;\;\overset{a.s.}{\to}\;\; 0,\;\;\;\;\;\;\frac{v_n(x)}{\tilde\sigma_n(z)}\;\;\overset{d}{\to}\;\;\mathcal N(0,1),
\end{align*}
where $z=L_{K^{1/2}}x$.
Furthermore, if $\{\varepsilon_i\}_{i=1}^n$ are Normally distributed, then $v_n(x)/\tilde\sigma_n(z)$ is exactly distributed as $\mathcal N(0,1)$, and
$\sup_{x\in\mX}|b_n(x)|/\hat\sigma_n^{\iso}(z)\to 0$ almost surely, and both are true without assuming (B3).
\label{thm:isokrr-limiting}
\end{theorem}

The following corollary is then immediate, noting that $\tilde\sigma_n(z)\leq\hat\sigma_n^\iso(z)$ almost surely.
\begin{corollary}
Impose conditions in Theorem \ref{thm:isokrr-limiting}. Then
$$
\liminf_{n\to\infty}\inf_{\beta_0\in\Theta_2(\mH)}\Pr[\eta_0(x)\in\varphi^\iso(x)]\;\;\geq\;\; 1-\alpha,\;\;\;\;\;\;\forall x\in\mX.
$$
Furthermore, if $\{\varepsilon_i\}_{i=1}^n$ are Normally distributed, then the following strengthened result holds and it holds without assuming (B3):
$$
\liminf_{n\to\infty}\inf_{P_X^n}\inf_{x\in\supp(P_X^n)}\inf_{\beta_0\in\Theta_2(\mH)}\Pr[\eta_0(x)\in\varphi^\iso(x)] \;\;\geq\;\; 1-\alpha.
$$
\label{cor:isokrr-ci}
\end{corollary}

\subsection{Rates of convergence}

\begin{theorem}
Suppose Assumptions (A1)-(A3) and (B1)-(B3) hold, with $\kappa<1/2$. Then as $n\to\infty$, 
\begin{equation}
\sup_{P_X^n}\sup_{\beta_0\in\Theta_2(\mH)}\mathbb E_{\eta_0}\left[\int\big|\varphi^\iso(x)\big|^2\ud P_X^n(x)\right]\lesssim n^{-\frac{2r-1}{2r}+\delta}
\label{eq:simple-validity}
\end{equation}
for any $\delta>0$.
\label{thm:isokrr-rates}
\end{theorem}


\subsection{Minimax optimality}

\begin{theorem}
Fix arbitrary $\alpha\in(0,1/4)$, $r>1/2$ and $\kappa=0$. There exists a distribution $P_X$ satisfying Assumptions (A1), (A3), (B1), (B2), and (B3), such that for any $\varphi$ satisfying Eq.~(\ref{eq:honest-ci})
with respect to the constructed $P_X$ and $\Theta_2(\mH)$, it holds that
$$
\sup_{\beta_0\in\Theta_2(\mH)} \mathbb E_{\beta_0}\left[\int\big|\varphi(x)\big|^2\ud P_X(x)\right] \gtrsim n^{-\frac{2r-1}{2r}}.
$$
\label{thm:isokrr-lb}
\end{theorem}

\begin{remark}
While not directly related, the problem instances constructed in the lower bound proof of Theorem \ref{thm:isokrr-lb} would also satisfy the spectral gap assumption (C1) to be introduced in later sections.
\end{remark}

We remark that, in the lower bound construction we fix the distribution of $X$ that satisfies all assumptions, including Assumption (B3), with $\kappa=0$.
Because this condition is stronger than Assumption (B3) for any $\kappa>0$, our lower bound implies lower bounds on relaxed assumptions of $\kappa>0$ as well.
Furthermore, because $P_X$ is fixed, an honest confidence interval can be considered to know $P_X$ in advance: our lower bounds show that even in this case,
an $\Omega(n^{-(2r-1)/2r})$ lower bound applies because the honest confidence interval does not know $\beta_0\in\Theta_2(\mH)$.

Theorem \ref{thm:isokrr-lb} is proved using similar constructions in \citep{cai2012minimax}, while borrowing ideas from \citep{cai2006adaptive} 
using two-point Le-Cam type arguments using Pinsker's inequality to prove a lower bound on the mean-squared widths of honest confidence intervals
that are significantly slower than optimal rates of convergence for functional regression estimators under the mean-squared error criterion.
It shows that, while known estimators in the functional regression setting attain $O(n^{-2r/(2r+1)}$ rate for the mean-squared error \citep{cai2012minimax},
the mean-squared widths of \emph{any} honest confidence interval must scale as $\Omega(n^{-(2r-1)/2r})$, effectively losing one-half degree of smoothness in the spectral decay rate of $T$.

\section{Under-smoothed anisotropic kernel regression}

We consider anisotropic extensions to the classical (isotropic) kernel ridge regression, by considering the following estimator
$$
\hat\theta_n^\ani = (T_n+\Lambda_n)^{\dagger}\frac{1}{n}\sum_{i=1}^n Y_i' Z_i'
$$
where $\Lambda_n$ is a positive-semidefinite, self-adjoint covariance operator computed on data that is not necessarily (multiples) of the identity operator,
and $C^\dagger$ is the Moore-Penrose pseudo-inverse of a not necessarily invertible linear operator $C$ (see, e.g.~\citep{fongi2023moore} and references therein).
Because $\Lambda_n$ must be entirely calculated from data, we consider the following form of $\Lambda_n$, rendering it a finite-rank operator:
$$
\Lambda_n = \sum_{j=1}^n \lambda_{nj}\hat\psi_j\otimes\hat\psi_j
$$
where $\{\lambda_{nj}\}_{j=1}^n\geq 0$ depends only on $\{Z_i'\}_{i=1}^n$ but not $Y_i'$.
With this choice, $\hat\theta_n^\ani$ can be written as
$$
\hat\theta_n^\ani = \frac{1}{n}\sum_{i=1}^n\sum_{j=1}^n \frac{Y_i'\langle Z_i',\hat\psi_j\rangle}{\hat s_j+\lambda_{nj}}\hat\psi_j.
$$

For the purpose of this section, we adopt the choice of 
$$
\lambda_{nj} =  \sqrt{\hat s_j\lambda_n}.
$$
With the estimator $\hat\eta^\ani_n(\cdot)$ such that $\hat\eta^\ani_n(x)=\hat\gamma_n+\int_I (L_{K^{1/2}}x)(t)\hat\beta^\ani_n(t)\ud t$  where $\hat\beta_n^\ani=L_{K^{1/2}}\hat\theta_n^\ani$, confidence intervals are constructed as
$$
\varphi^\ani(x) := [\hat\eta^\ani_n(x)\;\pm\;\tau_{\alpha/2}\hat\sigma_n^\ani(L_{K^{1/2}}x)]
$$
where $\tau_{\alpha/2}$ is the $(1-\alpha/2)$ quantile of the standard Normal distribution, and $\hat\sigma_n^\ani(\cdot)$ 
$$
\hat\sigma_n^\ani(z) = \frac{\sigma_\varepsilon}{\sqrt{n}}\sqrt{\langle z, (T_n+\Lambda_n)^{\dagger}T_n(T_n+\Lambda_n)^\dagger z\rangle}.
$$

\subsection{Honesty of confidence intervals}

\begin{theorem}
Suppose Assumptions (A1), (A2), (A3), (B1), (B2) hold. Suppose also that $\lambda_n$ satisfies Eq.~(\ref{eq:lambdan-1}).If 
\begin{equation}
\frac{\frac{1}{n}\sum_{i=1}^n |g(Z_1,Z_i)|^3}{\sqrt{n}(\frac{1}{n}\sum_{i=1}^n g(Z_1,Z_i)^2)^{3/2}}\;\;\overset{p}{\to}\;\; 0
\label{eq:anikrr-lyapnov}
\end{equation}
where $g(z,Z_i)=\langle z,(T_n+\Lambda_n)^{\dagger}Z_i\rangle$ and $Z_i=L_{K^{1/2}}X_i$, then it holds that
\begin{equation}
[\hat\sigma_n^\ani(Z_1)]^{-1}(\hat\eta_n^\ani(X_1)-\eta_0(X_1)) \;\;\overset{d}{\to}\;\;\mathcal N(0,1),
\label{eq:anikrr-clt}
\end{equation}
Furthermore, if $\{\varepsilon_i\}_{i=1}^n$ are Normally distributed, then Eq.~(\ref{eq:anikrr-clt}) holds for \emph{uniformly} for all $X_1,\cdots,X_n$ without the need to satisfy Eq.~(\ref{eq:anikrr-lyapnov}).
\label{thm:anikrr-limiting}
\end{theorem}

\begin{corollary}
Impose the same conditions in Theorem \ref{thm:anikrr-limiting}. Then 
$$
\liminf_{n\to\infty}\inf_{\beta_0\in\Theta_2(\mH)}\Pr[\eta_0(X_1)\in\varphi^\ani(X_1)] \;\; = \;\; 1-\alpha.
$$
where $X_1$ is the first input function in the sample $\{(X_i,Y_i)\}_{i=1}^n$ on which estimation $\hat\eta_n^{\ani}$ and confidence intervals $\varphi^\ani$.
Furthermore, if $\{\varepsilon_i\}_{i=1}^n$ are Normally distributed, the above result can be strengthened to
$$
\liminf_{n\to\infty}\min_{1\leq i\leq n}\inf_{\beta_0\in\Theta_2(\mH)}\Pr[\eta_0(X_i)\in\varphi^\ani(X_i)] \;\; = \;\; 1-\alpha.
$$
without the need to satisfy Eq.~(\ref{eq:anikrr-lyapnov}).
\label{cor:anikrr-ci}
\end{corollary}

\begin{remark}
In both statements of Theorem \ref{thm:anikrr-limiting} and Corollary \ref{cor:anikrr-ci}, we use the notation of $X_1$ and $Z_1$ without loss of generality.
Because $X_1,\cdots,X_n$ are i.i.d., $X_1,Z_1$ can be interpreted as an arbitrary input function in the sample $\{X_i,Y_i\}_{i=1}^n$.
\end{remark}

\subsection{Rates of convergence}

\begin{theorem}
Impose the same conditions in Theorem \ref{thm:anikrr-limiting}. Then the following holds almost surely as $n\to\infty$:
$$
\frac{1}{n}\sum_{i=1}^n \big|\varphi^\ani(X_i)\big|^2 \lesssim \sum_{j=1}^n \frac{\sqrt{\hat s_j}}{\sqrt{\hat s_j}+\sqrt{\lambda_n}}.
$$
where $X_1,\cdots,X_n$ are input functions on which $\hat\eta_n^\ani$ and $\varphi^\ani$ are estimated.
Furthermore, if $\hat s_j\lesssim j^{-2r}$ holds, then the right-hand side of the above inequality is upper bounded by $O(n^{-\frac{2r-1}{2r}+\delta})$
for any $\delta>0$.
\label{thm:anikrr-rates}
\end{theorem}

\subsection{Discussion of limitations}

{\color{red} Discuss the limitations and mention truncation is needed.}

\section{Truncated anisotropic kernel regression}\label{sec:anitr}

In this section we consider \emph{truncated} anisotropic kernel regression estimators, by setting $\lambda_{nj}\equiv +\infty$
for all $j> J(n)$ for a certain truncation level $J(n)\in[n]$ that depends on $n$.
This effectively mutes all eigenfunctions $\{\hat\psi_j\}_{j>J(n)}$, and the resulting estimator $\hat\theta_n$ admits the following spectral expression
\begin{equation}
\hat\theta_n^\anitr = \frac{1}{n}\sum_{i=1}^n\sum_{j=1}^{J(n)}\frac{Y_i'\langle Z_i',\hat\psi_j\rangle}{\hat s_j+\lambda_{nj}}\hat\psi_j.
\label{eq:anitr}
\end{equation}
For a compressed operator notation, define finite-rank operators $T_J$ and $\Lambda_J$ as
$$
T_J := \sum_{j=1}^{J(n)}\hat s_j\hat\psi_j\otimes\hat\psi_j,\;\;\;\;\;\; \Lambda_J := \sum_{j=1}^{J(n)}\lambda_{nj}\hat\psi_j\otimes\hat\psi_j.
$$
Then $\hat\theta_n^\anitr$ can be represented as
$$
\hat\theta_n^\anitr = (T_J+\Lambda_J)^{\dagger}\frac{1}{n}\sum_{i=1}^n Y_i'Z_i',
$$
where $C^\dagger$ is the Moore-Penrose pseudo-inverse of a not necessarily invertible linear operator $C$.

Following the previous section, one option for the choice of the anisotropic Ridge penalty parameters is $\lambda_{nj}=\lambda_n\wedge\sqrt{\hat s_j\lambda_n}$ with $\lambda_n$ satisfying Eq.~(\ref{eq:lambdan-1}),
and truncate at the level of $J(n)\leq n$.
Let $\{\omega_n\}_{n=1}^{\infty}$ be a sequence of slowly increasing positive constants such that
\begin{equation}
\omega_n\to\infty,\;\;\;\;\;\;\omega_n n^{-\delta}\to 0
\label{eq:omega}
\end{equation}
as $n\to\infty$, for any $\delta>0$.
The estimator is then $\hat\eta_n^\anitr(\cdot)=\hat\gamma_n + \langle\cdot,\hat\beta_n^\anitr\rangle$ where $\hat\beta_n^\anitr = L_{K^{1/2}}\hat\theta_n^\anitr$,
and $\langle f,g\rangle = \int_I f(t)g(t)\ud t$.
Afterwards, for a test function $x\in\supp(P_X)$, a confidence interval for $\eta_0(x)$ is constructed as
$$
\varphi^\anitr(x) = [\hat\eta_n^\anitr(x)\;\pm\; \tau_{\alpha}\hat\sigma_n^\anitr(L_{K^{1/2}}x) + \omega_n\|\hat\Pi_J^\perp (L_{K^{1/2}}x)\|_2],
$$
where $\tau_{\alpha}$ is the $(1-\alpha/2)$ quantile of the standard Normal distribution, $\hat\Pi_J^\perp z=(I-\sum_{j=1}^{J(n)}\hat\psi_j\otimes\hat\psi_j)z$, and $\hat\sigma_n^\anitr$ is defined as
$$
\hat\sigma_n^\anitr(z) = \frac{\sigma_\varepsilon}{\sqrt{n}}\sqrt{\langle z, (T_J+\Lambda_J)^\dagger T_J(T_J+\Lambda_J)^\dagger z\rangle}
$$


\subsection{Spectrum stability of the sample covariance operator}

In this section, we establish rigorously the connections between the (leading) spectrum of $T$ and $T_n$,
with cutoff thresholds set at levels that can be as deep as $o(\sqrt{n})$. Such discussion is important for the analysis of truncated anisotropic kernel regression estimators, for two primary reasons:
\begin{enumerate}
\item To establish asymptotic Normality for the widest class of testing functions $x\in\supp(P_X)$ possible, we must use the central limit theorem for independent but not identically distributed random variables
and verify classical conditions such as the Lyapnov condition or the Lindeberg condition. This involves understanding the distributions of ``influence functions'' in the variance term, which must be connected to
their population variants in order to avoid statistical correlation between $Z_i$ and $T_n$;
\item The condition (D1) to be introduced in the next section, unlike the condition (A2) that $\int_I \theta_0(t)^2\ud t\leq 1$, is \emph{not} basis invariant. This means that on a different orthonormal basis,
such as the eigenfunctions $\{\hat\psi_j\}_{j=1}^n$ of the sample covariance operator $T_n$, the sum $\sum_{j\geq 1}|\langle Z_j,\hat\psi_j\rangle|$ may not be bounded by $B$,
which leads to sub-optimal convergence rates if such sums indeed scale polynomially with $n$.
Therefore, we must rigorously compare the spectrum $\{\psi_j\}$ and $\{\hat\psi_j\}$ of the population and sample covariance operators,
to ensure that $\sum_{j\geq 1}|\langle Z_j,\hat\psi_j\rangle|$ remains bounded so that the bias of kernel Ridge regression can be effectively controlled.
\end{enumerate}

To establish such spectrum stability, we need the following assumption which lower bounds spectral gaps between consecutive eigenvalues on the population covariance operator:
\begin{enumerate}
\item[(C1)] There exists a numerical constant $c_g>0$ such that for all $j\geq 1$ \pd{need the gap condition for all $1 \leq j \leq \lceil C_J J(n) \rceil$, that is slightly above the truncation index}, $s_{j+1}-s_j\geq c_g j^{-2r-1}$, where $\{s_j\}_{j=1}^{\infty}$ are eigenvalues (listed in descending order) of the population covariance operator $T=\mathbb E_{P_X}[L_{K^{1/2}}X\otimes L_{K^{1/2}}X]$.
\end{enumerate}

We further impose the following condition on the estimator's choice of truncation levels.
\begin{enumerate}
\item[(C2)] There exists a constant $q\in(0,1/2)$ and $0<c_q\leq C_q<+\infty$ such that $c_q n^q\leq J(n)\leq C_q n^q$ for all $n\geq 3$.
\end{enumerate}

Recall that the population covariance operator $T=\mathbb E[Z\otimes Z]$ and the sample covariance operator $T_n=\frac{1}{n}\sum_{i=1}^nZ_i\otimes Z_i$ admit spectral decomposition of
$T=\sum_{j=1}^{\infty}s_j\psi_j\otimes\psi_j$ and $T_n=\sum_{j=1}^n\hat s_j\hat\psi_j\otimes\hat\psi_j$, where $\{\psi_j\}_{j=1}^{\infty}$ is an orthonormal basis of $L_2(I)$ and $\{\hat\psi_j\}_{j=1}^n$ are orthonormal in $L_2(I)$.
\begin{lemma}
Suppose Assumptions (A3), (B1), (B2), (B3), (C1) and (C2) hold, with parameters $\kappa<1/4$ and $q\in(0,1/2)$.
Fix an arbitrary small constant $c\geq 2$.
Then there exist constants $C_J,L_J<+\infty$ depending only on $c$ and other parameters in assumptions, such that for every $n\geq 3$, 
with probability $1-n^{-c}$ the following holds: there exists an injection $\varpi:[\lceil C_J J(n)\rceil]\to \mathbb N$ which is a surjection on $[J(n)]$,
\footnote{This means that 1) for every $j,j'\in[\lceil C_JJ(n)\rceil]$, $j\neq j'$ implies $\varpi(j)\neq\varpi(j')$, and 2) for every $\ell\in[J(n)]$, there exists $j\in[\lceil C_JJ(n)\rceil]$ such that $\varpi(j)=\ell$.}
 such that for every $j\in[\lceil C_J n\rceil]$, the following hold:
\begin{enumerate}
\item $j/C_J\leq \varpi(j)\leq C_J j$;
\item $|\hat s_{\varpi(j)}/s_j-1|\leq L_J j\log^3 n/\sqrt{n}$;
\item For every $k\in\mathbb N$, $k\neq j$, $|\langle\hat\psi_{\varpi(j)},\psi_k\rangle|\leq L_J\Delta_\psi(j,k)$, where
$$
\Delta_\psi(j,k) = \frac{\log n\log(3j)\log(3k)}{\sqrt{n}}\times \left\{\begin{array}{ll} k^r/j^r,& \text{if }k< j/C_J;\\ j^r/k^r,& \text{if } k> C_Jj;\\ \sqrt{jk}/|j-k|;& \text{if $j/C_J \leq k \leq C_Jj$ and $j\neq k$.}\end{array}\right.
$$
\end{enumerate}
Furthermore, if $q<1/4$ then $\varpi$ can be chosen as the bijection $\varpi(j)=j$ for all $j\leq \lceil C_JJ(n)\rceil$.
\label{lem:stability}
\end{lemma}

%\begin{enumerate}
%\item For every $j\leq J(n)$, there exists $j'\leq C_J j\leq C_JJ(n)$ such that $\|\hat\psi_j-\psi_{j'}\|_2\leq L_J j\sqrt{\log n}/\sqrt{n}$ and $|\hat s_j/s_{j'}-1|\leq L_J j\sqrt{\log n}/\sqrt{n}$;
%additionally, for every $k\in[C_J J(n)]\backslash\{j'\}$, $|\langle\hat\psi_j,\psi_k\rangle|\leq L_J\Delta_\psi(j',k)$;
%\item For every $j'\leq J(n)/C_J$, there exists $j\leq C_Jj'\leq J(n)$ such that $\|\hat\psi_j-\psi_{j'}\|_2\leq L_J j'\sqrt{\log n}/\sqrt{n}$ and $|\hat s_j/s_{j'}-1|\leq L_J j'\sqrt{\log n}/\sqrt{n}$;
%additionally, for every $k\in[C_J J(n)]\backslash\{j'\}$, $|\langle\hat\psi_j,\psi_k\rangle|\leq L_J\Delta_\psi(j',k)$.
%\end{enumerate}


{\color{red} Mention why this can't be done using Davis Kahan, citing existing fPCA papers.}

{\color{red} Mention also that the $q<1/4$ bijection result is for information only and \emph{not} necessary for our results.}


The following lemma, as a consequence of Lemma \ref{lem:stability} and Assumptions (B2) and (B3) regulating the behaviors of $\langle Z,\psi_j\rangle$,
upper bounds the discrepancy between coefficients on $\{\psi_j\}$ and $\{\hat\psi_j\}$ of a test function $z=L_{K^{1/2}}x$.
\begin{lemma}
Impose the same set of conditions in Lemma \ref{lem:stability}.
Let $J=J(n)\asymp n^q$ and $\und J=J(n)/C_J$.
For each $M\in\{2,4\}$, conditioned on results in Lemma \ref{lem:stability},
it holds that
$$
\sup_{\ell\leq J}\sup_{z\in\supp(P_Z^n)}\big|\langle z,\hat\psi_\ell\rangle^M - \langle z,\psi_j\rangle^M\big|\; \lesssim\; n^{M\kappa}\times \frac{J^{1-rM}\log^4 n}{\sqrt{n}},
$$
where $j$ is the unique integer satisfying $\varpi(j)=\ell$, whose existence is guaranteed by Lemma \ref{lem:stability}.
Furthermore, 
$$
\sup_{\ell\leq J}\bigg|\mathbb E_{Z\sim P_Z^n}\left[\langle Z,\hat\psi_\ell\rangle^M\right] -\mathbb E_{Z\sim P_Z^n}\left[\langle Z,\psi_j\rangle^M\right]\bigg|
\;\lesssim\;\frac{J^{1-rM}\log^4 n}{\sqrt{n}}.
$$
\label{lem:zpsi-stability}
\end{lemma}

\subsection{Honesty of confidence intervals}

The following theorem establishes the limiting distribution of $\hat\eta_n^\anitr$.
\begin{theorem}
Impose all assumptions and conditions in Lemma \ref{lem:stability}. Suppose $\Lambda_n$ is constructed as $\lambda_{nj}=\sqrt{\hat s_j\lambda_n}$ for all $j\leq J(n)$,
where $\lambda_n$ satisfies Eq.~(\ref{eq:lambdan-1}). Suppose also that Assumptions (A1), (A2) hold, and $\kappa<(1-q)/2$. Then 
$$
\lim_{n\to\infty} \sup_{P_X^n}\sup_{x\in\supp(P_X^n)}\sup_{\beta_0\in\Theta_2(\mH)}\sup_{t\in\mathbb R}\left|\Pr\left[\frac{\hat\eta_n^{\anitr}(x)-\eta_0(x)+\langle\hat\Pi_J^\perp z, \theta_0\rangle}{\hat\sigma_n^\anitr(L_{K^{1/2}})}\leq t\right] - \Phi(t)\right| = 0,
$$
where $z=L_{K^{1/2}}x$, $\theta_0=L_{K^{-1/2}}\beta_0$, $\hat\Pi_J^\perp = I - \sum_{j=1}^{J(n)}\hat\psi_j\otimes\hat\psi_j$, and $\Phi(\cdot)$ is the CDF of the standard Normal distribution.
The probability is over the randomness of both $\{X_i\}_{i=1}^n$ and $\{\varepsilon_i\}_{i=1}^n$. 

{\color{red}
Furthermore, if $\frac{1}{2r}<q<\frac{1}{2}$ then for every $\epsilon>0$, there exist measurable sets $\mX_1\subseteq\supp(P_X^1),\mX_2\subseteq\supp(P_X^2),\cdots$ satisfying $P_X^1(\mX_1),P_X^2(\mX_2),\cdots\geq 1-\epsilon$,
such that
$$
\lim_{n\to\infty} \sup_{P_X^n}\sup_{x\in\mX_n}\sup_{\beta_0\in\Theta_2(\mH)}\sup_{t\in\mathbb R}\left|\Pr\left[\frac{\hat\eta_n^{\anitr}(x)-\eta_0(x)}{\hat\sigma_n^\anitr(L_{K^{1/2}}z)}\leq t\right] - \Phi(t)\right| = 0.
$$}
\label{thm:anitr-limiting}
\end{theorem}

{\color{red} YW: Furthermore part is still troublesome, arising from a collection of cross terms not handled properly in a technical lemma.}

\begin{remark}
Taking $q$ to be close to $1/2$, the condition $1/2r<q<1/2$ amounts to $r>1$.
\end{remark}

Note that 
\begin{equation}
\big|t_n(z)\big| = \big|\langle z, \hat\Pi_J^\perp\theta_0\rangle\big|\leq \|\hat\Pi_J^\perp z\|_2\cdot \|\theta_0\|_2 \leq \omega_n\|\hat\Pi_J^\perp z\|_2,
\label{eq:anitr-limiting-impl}
\end{equation}
where the first inequality holds from Cauchy-Schwarz inequality and the fact that $\hat\Pi_J^\perp$ is a self-adjoint, projection operator, and the 
last inequality holds for sufficiently large $n$ because $\|\theta_0\|_2\leq 1$ and $\omega_n\to\infty$ as $n\to\infty$.

\begin{corollary}
Impose the same set of assumptions and conditions in Theorem \ref{thm:anitr-limiting}. Then
$$
\liminf_{n\to\infty}\inf_{P_X^n}\inf_{x\in\supp(P_X^n)}\inf_{\beta_0\in\Theta_2(\mH)}\Pr[\eta_0(x)\in\varphi^\anitr(x)] \;\;\geq\;\; 1-\alpha.
$$
Furthermore, if $\frac{1}{2r}<q<\frac{1}{2}$, then the following holds for the choice of $\omega_n\equiv 0$:
for any $\epsilon>0$, there exist measurable sets $\mX_1\subseteq\supp(P_X^1),\mX_2\subseteq\supp(P_X^2),\cdots$ satisfying $P_X^1(\mX_1),P_X^2(\mX_2),\cdots\geq 1-\epsilon$,
such that for any sequence $x_1\in\mX_1,x_2\in\mX_2,\cdots$, 
$$
\lim_{n\to\infty}\inf_{\beta_0\in\Theta_2(\mH)} \Pr[\eta_0(x)\in\varphi^\anitr(x)] \;\; = \;\; 1-\alpha.
$$
\label{cor:anitr}
\end{corollary}

\subsection{Rates of convergence}

\begin{theorem}
Impose all assumptions and conditions in Lemma \ref{lem:stability}. Suppose $\Lambda_n$ is constructed as $\lambda_{nj}=\sqrt{\hat s_j\lambda_n}$ for all $j\leq J(n)$,
where $\lambda_n$ satisfies Eq.~(\ref{eq:lambdan-1}). Then as $n\to\infty$, 
$$
\sup_{P_X^n}\sup_{\beta_0\in\Theta_2(\mH)}\mathbb E\left[\int \big|\varphi^\anitr(x)\big|^2\ud P_X^n(x) \right] \lesssim n^{-\frac{2r-1}{2r}+\delta} + n^{2q-\frac{3}{2}+\delta} + n^{-q(2r-1)+\delta} + n^{1-c}.
$$
for any $\delta>0$, where $c\geq 2$ is the parameter in the failure probability in Lemma \ref{lem:stability}.
\label{thm:anitr-rates}
\end{theorem}

\begin{remark}
If $\frac{1}{2r}<q<\frac{1}{4r}+\frac{1}{4}$ then the right-hand side of the above inequality is asymptotically on the order of $O(n^{-(2r-1)/2r+\delta})$ for any $\delta>0$.
\end{remark}


%Let $\mH\subseteq L_2(I)$ be a sufficiently smooth subset of all square-integrable functions on $L_2(I)$. For example, $\mH$ could be the Sobolev class with smoothness level $\rho>1/2$.
%Let $X_1,\cdots,X_n$ be i.i.d.~supported on $\mH$, and $\varepsilon_1,\cdots,\varepsilon_n\overset{i.i.d.}{\sim}\mathcal N(0,1)$. Is it true that, with probability $1-n^{-1}$,
%$$
%\sup_{\psi\in\mH}\frac{\frac{1}{n}\sum_{i=1}^n \langle X_i,\psi\rangle^2-\mathbb E[\langle X,\psi\rangle^2]}{\sqrt{\mathbb E[\langle X,\psi\rangle^4]}} \lesssim \gamma(\mH)\sqrt{\frac{\log n}{n}}?
%$$
%And
%$$
%\sup_{\psi\in\mH}\frac{\frac{1}{n}\sum_{i=1}^n \varepsilon_i\langle X_i,\psi\rangle}{\sqrt{\mathbb E[\langle X,\psi\rangle^2]}}\lesssim \gamma(\mH)\sqrt{\frac{\log n}{n}}?
%$$
%Is there a reference for both inequalities? What is the form of $\gamma(\mH)$ and how does it depend on $\rho$? If $1/\sqrt{n}$ is not possible, what is the right rate?

\section{More adaptivity in anisotropic regularization weights}\label{sec:adptr}

\begin{enumerate}
\item[(C1)] There exists a numerical constant $B_1<+\infty$ such that $\sum_{j=1}^{\infty}|\langle \theta_0, \psi_j\rangle| \leq B_1$, where $\{\psi_j\}_{j=1}^{\infty}$ are eigenfunctions of the population covariance operator $T$,
and $\theta_0=L_{K^{-1/2}}\beta_0$ for some $\beta_0\in\mH$.
\end{enumerate}
We use $\Theta_1(\mH)\subseteq\Theta_2(\mH)$ to denote all unknown linear functionals in $\Theta_2(\mH)$ that also satisfies Assumption (D1), for a pre-determined constant upper bound $B_1$.
This assumption is also implicitly imposed, albeit in a different form, in the works of \cite{liu2025scalable} to construct pointwise confidence intervals for classical kernel estimators.

{\color{red}Maybe some remarks on Fourier Basis, or alignment arising from Gaussian processes.}

Let $\mu_n>0$ be a sequence of penalty parameters that decay slightly slower than $\lambda_n$, satisfying
\begin{equation}
\mu_n n^{2r/(2r+1)}\to 0,\;\;\;\;\;\;\mu_n n^{2r/(2r+1)+\delta}\to\infty
\label{eq:mun}
\end{equation}
for all $\delta>0$.
We consider $\hat\theta_n^\adptr$ as solution to Eq.~(\ref{eq:anitr}), with $\Lambda_n=\sum_{j=1}^{J(n)}\lambda_{nj}(z)\hat\psi_j\otimes\hat\psi_j$ being defined as
\begin{equation}
\lambda_{nj}(z) = \left\{\begin{array}{ll}\hat s_j\sqrt{\mu_n}/|\langle z,\hat\psi_j\rangle|,&\text{if } \langle z,\hat\psi_j\rangle^2\geq 2\hat s_j;\\  \sqrt{\hat s_j\mu_n},& \text{otherwise};\end{array}\right.
\label{eq:lambda-adp}
\end{equation}
and truncation occurs at $J(n)\leq n$.
Note that in this estimator, the choice of $\Lambda_n$ and the penalization parameters $\lambda_{nj}$ are \emph{adaptive} to the testing function $z=L_{K^{1/2}}x$,
which is different from the previous estimators that do not depend on $z$.

We remark that the constant $2$ in the first prong of the definition of Eq.~(\ref{eq:lambda-adp}) can be changed to any constant without affecting the analysis of this estimator.
This also implies that, if Assumption (B3) holds with $\kappa=0$, then there is no need to do have anisotropic regularization parameters depending on
the test function $z$, and one can instead set simply $\lambda_{nj}= \mu_n\wedge\sqrt{\hat s_j\mu_n}$.

For a testing function $z$ and point estimate $\hat\eta_n^\adptr(x)=\hat\gamma_n+\langle L_{K^{1/2}}z,\hat\theta_n^\adptr\rangle$, a confidence interval of $\eta_0(x)$ is constructed as
$$
\varphi^\adptr(x) = [\hat\eta_n^\adptr(x)\;\pm\;\tau_{\alpha}\hat\sigma_n^\adptr(L_{K^{1/2}}x)+\omega_n(\sqrt{\mu_n}+\|\hat\Pi_J^\perp(L_{K^{1/2}}x)\|_2)],
$$
where $\tau_{\alpha}$ is the $(1-\alpha/2)$ quantile of the standard Normal distribution, and $\hat\sigma_n^\adptr$ is defined as
$$
\hat\sigma_n^\adptr(z) =  \frac{\sigma_\varepsilon}{\sqrt{n}}\sqrt{\langle z, (T_J+\Lambda_J)^\dagger T_J(T_J+\Lambda_J)z\rangle} = \frac{\sigma_\varepsilon}{\sqrt{n}}\sqrt{\sum_{j=1}^{J(n)}\frac{\langle z,\hat\psi_j\rangle^2\hat s_j}{(\hat s_j+\lambda_{nj})^2}}.
$$

\subsection{Honesty of confidence intervals}

\begin{theorem}
Impose all assumptions and conditions in Lemma \ref{lem:stability}. Suppose $\Lambda_n$ is constructed in Eq.~(\ref{eq:lambda-adp}) for all $j\leq J(n)$,
where $\mu_n$ satisfies Eq.~(\ref{eq:mun}) and $q$ in Assumption (C2) satisfies $q<1/4$. Suppose also that Assumptions (A1), (A2), (D1) hold, and $\kappa<(1-q)/2$. Then 
$\hat\eta_n^\adptr(x)-\eta_0^\adptr(x) \pd{\eta_0(x)}+\langle\hat\Pi_J^\perp L_{K^{1/2}}x,\theta_0\rangle = b_n(L_{K^{1/2}}x)+v_n(L_{K^{1/2}}x)$ for every $n$ and $x\in\supp(P_X^n)$, where
$$
\limsup_{n\to\infty}\sup_{P_Z^n}\sup_{z\in\supp(P_Z^n)}\sup_{\beta_0\in\Theta_1(\mH)} \frac{|b_n(z)|}{\sqrt{\mu_n}} < +\infty\;\;\;\;\;\;a.s.,
$$
and
$$
\lim_{n\to\infty} \sup_{P_Z^n}\sup_{z\in\supp(P_Z^n)}\sup_{\beta_0\in\Theta_1(\mH)}\sup_{t\in\mathbb R}\left|\Pr\left[\frac{v_n(z)}{\hat\sigma_n^\adptr(z)}\leq t\right] - \Phi(t)\right| = 0,
$$
where $\theta_0=L_{K^{-1/2}}\beta_0$, $\hat\Pi_J^\perp = I - \sum_{j=1}^{J(n)}\hat\psi_j\otimes\hat\psi_j$, and $\Phi(\cdot)$ is the CDF of the standard Normal distribution.
The probability is over the randomness of both $\{X_i\}_{i=1}^n$ and $\{\varepsilon_i\}_{i=1}^n$. 

Furthermore, if $\frac{2r}{(2r-1)(2r+1)}<q<\frac{1}{4}$, then for every $\epsilon>0$, there exist measurable sets $\mX_1\subseteq\supp(P_X^1),\mX_2\subseteq\supp(P_X^2),\cdots$ satisfying $P_X^1(\mX_1),P_X^2(\mX_2),\cdots\geq 1-\epsilon$,
such that for any sequence $x_1\in\mX_1,x_2\in\mX_2,\cdots$, it holds that
$$
\lim_{n\to\infty} \sup_{\beta_0\in\Theta_1(\mH)}\sup_{t\in\mathbb R}\left|\Pr\left[\frac{\hat\eta_n^\adptr(x)-\eta_0(x)}{\hat\sigma_n^\adptr(L_{K^{1/2}}x)}\leq t\right] - \Phi(t)\right| = 0.
$$
%
%Furthermore, if $??$ then for every $\epsilon>0$, there exist measurable sets $\mX_1\subseteq\supp(P_X^1),\mX_2\subseteq\supp(P_X^2),\cdots$ satisfying $P_X^1(\mX_1),P_X^2(\mX_2),\cdots\geq 1-\epsilon$,
%such that
%$$
%\lim_{n\to\infty} \sup_{P_X^n}\sup_{x\in\mX_n}\sup_{\beta_0\in\Theta_1(\mH)}\sup_{t\in\mathbb R}\left|\Pr\left[\frac{\hat\eta_n^{\adptr}(x)-\eta_0(x)}{\hat\sigma_n^\adptr(L_{K^{1/2}}z)}\leq t\right] - \Phi(t)\right| = 0.
%$$
\label{thm:adptr-limiting}
\end{theorem}

\begin{corollary}
Impose the same set of assumptions and conditions in Theorem \ref{thm:adptr-limiting}. Then
$$
\liminf_{n\to\infty}\inf_{P_X^n}\inf_{x\in\supp(P_X^n)}\inf_{\beta_0\in\Theta_2(\mH)}\Pr[\eta_0(x)\in\varphi^\adptr(x)] \;\;\geq\;\; 1-\alpha.
$$
Furthermore, if $\frac{2r}{(2r-1)(2r+1)}<q<\frac{1}{4}$, then the following holds for the choice of $\omega_n\equiv 0$:
for any $\epsilon>0$, there exist measurable sets $\mX_1\subseteq\supp(P_X^1),\mX_2\subseteq\supp(P_X^2),\cdots$ satisfying $P_X^1(\mX_1),P_X^2(\mX_2),\cdots\geq 1-\epsilon$,
such that for any sequence $x_1\in\mX_1,x_2\in\mX_2,\cdots$, 
$$
\lim_{n\to\infty}\inf_{\beta_0\in\Theta_1(\mH)} \Pr[\eta_0(x)\in\varphi^\adptr(x)] \;\; = \;\; 1-\alpha.
$$
\label{cor:adptr}
\end{corollary}



\subsection{Rates of convergence}


\begin{theorem}
Impose all assumptions and conditions in Lemma \ref{lem:stability}. Suppose $\Lambda_n$ is constructed as in Eq.~(\ref{eq:lambda-adp}),
where $\mu_n$ satisfies Eq.~(\ref{eq:mun}). Then as $n\to\infty$, 
$$
\sup_{P_X^n}\sup_{\beta_0\in\Theta_1(\mH)}\mathbb E\left[\int \big|\varphi^\adptr(x)\big|^2\ud P_X^n(x) \right] \lesssim n^{-\frac{2r}{2r+1}+\delta} + n^{2q-\frac{3}{2}+\delta} + n^{-q(2r-1)+\delta} + n^{1-c}.
$$
for any $\delta>0$, where $c\geq 2$ is the parameter in the failure probability in Lemma \ref{lem:stability}.
\label{thm:adptr-rates}
\end{theorem}

\begin{remark}
For both Theorems \ref{thm:adptr-limiting} and \ref{thm:adptr-rates} to hold with an $n^{-2r/(2r+1)+\delta}$ mean-squared confidence interval width, $q$ must satisfy
$$
\frac{2r}{(2r-1)(2r+1)}\;<\; q\;<\; \frac{1}{2}.
$$
Provided that $r>(1+\sqrt{2})/2=1.207$, the valid range of $q$ is non-empty.
\end{remark}


The ``effective dimension'' $d$ at which $s_d\asymp \mu_n$, is $d\asymp n^{1/(2r+1)}$ when $\mu_n\approx n^{-2r/(2r+1)}$.
This means that the ``truncation level'' $J$ should be significantly larger than $d$, because $q>\frac{2r}{(2r-1)(2r+1)} > \frac{1}{2r+1}$.
This means that, in the anisotropic setting, \emph{both} truncation and Ridge smoothing are necessary to achieve optimal honest and adaptive confidence intervals.

\subsection{Minimax optimality of $\varphi^{\adptr}$}


\section{Proofs}

For brevity we assume $\mathbb E[x]=0$, $\gamma_0=0$ and $X_i'=X_i$ throughout all proofs, which is conventional in the analysis of linear functional regression in the literature \citep{cai2012minimax,yuan2010reproducing}.
For all \emph{lower bound} proofs, we also assume $L_{K^{1/2}}=L_{K^{-1/2}}$ is the identity operator to simplify the proofs.
For notational brevity, we many times drop the identity operator symbol $I$, such as writing $T_n+\lambda$ for $T_n+\lambda I$, whenever there is no ambiguity.



\subsection{Some technical lemmas for results in Sec.~\ref{sec:isotropic}}

We first establish a few technical lemmas which are important to our proof of later results.
\begin{lemma}
Suppose Assumptions (A1), (A3) and (B3) hold with $\kappa<\frac{1}{2}-\frac{1}{4r}$, and $\lambda_n$ satisfies Eq.~(\ref{eq:lambdan-1}). Then for any $\delta\in(0,1/2]$ there exists a constant $n_0\in\mathbb N$ depending only on $\delta$ and other parameters in assumptions,
 such that for all $n\geq n_0$, with probability at least $1-n^{-1}$,
$$
(1-\delta)(T+\lambda_n)\preceq T_n+\lambda_n \preceq (1+\delta)(T+\lambda_n)
$$
where $T_n=\frac{1}{n}\sum_{i=1}^n Z_i\otimes Z_i$, $T=\mathbb E[Z\otimes Z]$, $Z_i=L_{K^{1/2}}X_i$, $Z=L_{K^{1/2}}Z$, and for two self-adjoint operators $A,B$, $A\preceq B$ if $\langle x,Ax\rangle\leq \langle x, Bx\rangle$
for any square integrable on function $x:I\to\mathbb R$.
Furthermore, for any $\delta>0$, there exists a constant $C<+\infty$ depending only on parameters in the assumptions, such that for all $n\geq 3$, with probability $1-n^{-1}$ it holds that
$$
\left\|(T+\lambda_n)^{-1/2}(T-T_n)(T+\lambda_n)^{-1/2}\right\|_\op \leq C\times n^{-\frac{1-2\kappa-1/2r}{2}+\delta}.
$$
\label{lem:empirical-population-cov}
\end{lemma}

We remark that similar versions of Lemma \ref{lem:empirical-population-cov} have appeared, in different forms, in classical analysis of kernel ridge regression estimators and RKHS approaches for linear functional regression.
See e.g.~the upper bound on $\mathcal S_2(\lambda,z)$ in \citep{caponnetto2007optimal}, the second inequality in the proof of Theorem 5 in \citep{smale2007learning}, 
Theorem 2 of \citep{hsu2014random},
or Lemma 2 of \citep{cai2012minimax}.
Nevertheless, in our applications the choice of $\lambda_n$ is significantly smaller than conventional choices (for inference purposes), and when $\lambda_n = o(n^{-1})$ existing arguments cannot be directly applied,
without additional conditions such as Assumption (B3) for some $\kappa<1/2$. Therefore, we give the complete proof of Lemma \ref{lem:empirical-population-cov} below.

\begin{proof}[Proof of Lemma \ref{lem:empirical-population-cov}.]
We use Hoeffding's inequality for bounded self-adjoint operators, as developed and stated in \citep{minsker2017some}.
For every $i\in[n]$, define
$$
W_i := (T+\lambda_n)^{-1/2}Z_i = (T+\lambda_n)^{-1/2}L_{K^{1/2}}X_i.
$$
This gives rise to
\begin{align}
\frac{1}{n}\sum_{i=1}^n W_i\otimes W_i &= (T+\lambda_n )^{-1/2}T_n (T+\lambda_n )^{-1/2};\label{eq:epcov-0}\\
\mathbb E[W_i\otimes W_i] &= (T+\lambda_n)^{-1/2} T (T+\lambda_n )^{-1/2}, \;\;\;\;\;\;\forall i\in[n]\label{eq:epcov-0half}.
\end{align}

By Assumptions (B1) and (B3), it holds almost surely that 
\begin{align}
\|W_i\otimes W_i\|_\op&=\left\|\sum_{j,k=1}^{\infty}\frac{\langle Z_i,\psi_j\rangle\langle Z_i,\psi_k\rangle}{\sqrt{s_j+\lambda_n}\sqrt{s_k+\lambda_n}}\psi_j\otimes\psi_k\right\|_\op\nonumber\\
&\leq \sum_{j=1}^{\infty}\frac{\langle Z_i,\psi_j\rangle^2}{s_j+\lambda_n} \leq \sum_{j=1}^{\infty} \frac{C_\nu s_j n^{2\kappa}}{s_j+\lambda_n}  \lesssim n^{2\kappa+1/2r+\delta},
\label{eq:proof-epcov-1}
\end{align}
where the last inequality holds because $\lambda_n\gtrsim n^{-1-\delta}$ for any $\delta>0$. 

Let 
$$
A_n := \sum_{i=1}^n\mathbb E[(W_i\otimes W_i)^2] =n\times\sum_{j,k=1}^{\infty}\frac{\mathbb E[\langle Z,Z\rangle\langle Z,\psi_j\rangle\langle Z,\psi_k\rangle]}{(s_j+\lambda_n)(s_k+\lambda_n)}\psi_j\otimes \psi_k.
$$
It is easy to verify that $\|A_n\|_\op \geq n\mathbb E[\langle Z,\psi_1\rangle^4]/(s_1+\lambda_n)= \pd{\geq \ }  ns_1^2/(s_1+\lambda_n)\geq c_sn/2$ for sufficiently large $n$. 
On the other hand, Because $\langle Z,Z\rangle\leq 1$ almost surely thanks to Assumption (A3), it holds that
\begin{equation}
\tr(A_n) \leq n\times \sum_{j=1}^{\infty} \frac{s_j}{(s_j+\lambda_n)^2} \lesssim n^{(2r+1)/r},
\label{eq:proof-epcov-2}
\end{equation}
where the last inequality holds because $s_j\asymp j^{-2r}$ thanks to Assumption (A1), and $\lambda_n\gtrsim n^{-2}$.
Invoking \citep[Theorem 3.1]{minsker2017some}, for any $\delta>0$, it holds with probability $1-n^{-1}$ that
\begin{align*}
\left\|\frac{1}{n}\sum_{i=1}^n W_i\otimes W_i-\mathbb E[W_i\otimes W_i]\right\|_\op &\lesssim n^{-\frac{1-2\kappa-1/2r-\delta}{2}}\log n = o(1). 
\end{align*}
Subsequently, incorporating Eqs.~(\ref{eq:epcov-0},\ref{eq:epcov-0half}), the above inequality implies
$$
\left\|(T+\lambda_n)^{-1/2}(T-T_n)(T+\lambda_n)^{-1/2}\right\|_\op \lesssim n^{-\frac{1-2\kappa-1/2r}{2}+\delta} = o(1).
$$
This proves Lemma \ref{lem:empirical-population-cov}, by selecting $n_0$ sufficiently large such that the right-hand side of the above inequality is upper bounded by $\delta$ for any $\delta>0$.
\end{proof}

For the purpose of the next lemma, for any $z=L_{K^{1/2}}x$ for some $x\in\mathcal X$ and $Z_i=L_{K^{1/2}}X_i$, define
$$
g(z,Z_i) := \langle z, (T_n+\lambda_n)^{-1}Z_i\rangle,\;\;\;\;\;\;\bar g(z,Z_i) := \langle z, (T+\lambda_n)^{-1}Z_i\rangle.
$$
The following lemma derives probabilistic convergence between $g$, $\bar g$, and their expectations.
\begin{lemma}
Suppose the same set of conditions in Lemma \ref{lem:empirical-population-cov} hold, with Assumption (B3) being strengthened to $\kappa<\frac{1}{6}-\frac{1}{4r}$. 
Then it holds that
\begin{align*}
\sup_{z\in L_{K^{1/2}}\mX}\left|\frac{1}{n}\sum_{i=1}^n [g(z,Z_i)-\bar g(z,Z_i)]^2\right|\;\;&\overset{p}{\to}\;0;\\
\sup_{z\in L_{K^{1/2}}\mX}\left|\frac{1}{n}\sum_{i=1}^n \bar g(z,Z_i)^2-\mathbb E[\bar g(z,Z)^2]\right|\;\;&\overset{p}{\to}\; 0.
\end{align*}
\label{lem:gdiff}
\end{lemma}
\begin{proof}[Proof of Lemma \ref{lem:gdiff}.]
Throughout this proof, we write $w := (T+\lambda)^{-1/2} z$ and $W_i := (T_n+\lambda)^{-1/2}Z_i$ for notational simplicity, where $\lambda=\lambda_n$.
By Assumptions (B1) and (B3), for sufficiently large $n$ it holds uniformly for almost every $x\in\mX$ and $z=L_{K^{1/2}}x$ that
\begin{equation}
\|w\|_2^2 \leq \sum_{j=1}^{\infty} \frac{\langle z,\psi_j\rangle^2}{s_j+\lambda_n} \leq \sum_{j=1}^{\infty}\frac{C_\nu s_j n^{2\kappa}}{s_j+\lambda_n} \lesssim n^{2\kappa+1/2r+\delta}
\label{eq:proof-gdiff-0}
\end{equation}
for any $\delta>0$. 
%\pd{Need to be a little careful here as (B3) holds almost surely for $Z$ but in (22), the upper bound on $\langle z,\psi_j\rangle$ needs to hold for all $z$. This is fine, but just need to adapt wording.}

Define $\Delta_i := g(z,Z_i)-\bar g(z,Z_i)$. Then
\begin{align*}
\Delta_i &= \langle z,((T_n+\lambda )^{-1}-(T+\lambda )^{-1})Z_i\rangle\nonumber\\
&= \langle z, (T+\lambda )^{-1}((T+\lambda )(T_n+\lambda )^{-1}-I)Z_i\rangle\nonumber\\
&= \langle z, (T+\lambda )^{-1}(T_n-T)(T_n+\lambda )^{-1}Z_i\rangle.
\end{align*}
Subsequently,
\begin{align}
\frac{1}{n}\sum_{i=1}^n \Delta_i^2 &= \langle z, (T+\lambda)^{-1}(T_n-T)(T_n+\lambda)^{-1}T_n(T_n+\lambda)^{-1}(T_n-T)(T+\lambda)^{-1}z\rangle\nonumber\\
&= \langle w, (T+\lambda)^{-1/2}(T_n-T)(T_n+\lambda)^{-1}T_n(T_n+\lambda)^{-1}(T_n-T)(T+\lambda)^{-1/2}w\rangle\nonumber\\
&\leq \langle w, (T+\lambda)^{-1/2}(T_n-T)(T_n+\lambda)^{-1}(T_n-T)(T+\lambda)^{-1/2}w\rangle\nonumber\\
&\leq \|(T+\lambda)^{-1/2}(T_n-T)(T_n+\lambda)^{-1/2}\|_\op\|w\|_2^2\nonumber\\
&\overset{(a)}{=} O_P\left(n^{-\frac{1-2\kappa-1/2r}{2}+\delta}\right)\times O(n^{2\kappa+1/2r+\delta}) \;\;\overset{p}{\to} \;\;0,
\end{align}
where Eq.~(a) holds by invoking Lemma \ref{lem:empirical-population-cov} and incorporating Eq.~(\ref{eq:proof-gdiff-0}), {\color{magenta}and noting that
\begin{align*}
\|(T+\lambda)^{-1/2}&(T_n-T)(T_n+\lambda)^{-1/2}\|_\op \nonumber\\
&\leq \|(T+\lambda)^{-1/2}(T_n-T)(T+\lambda)^{-1/2}\|_\op \cdot \|(T+\lambda)^{1/2}(T_n+\lambda)^{-1/2}\|_\op,
\end{align*}
where the first term can be upper bounded directly using Lemma \ref{lem:empirical-population-cov}, and the second term can be upper bounded by noting that 
$$
\|(T+\lambda)^{-1/2}(T_n-T)(T_n+\lambda)^{-1/2}\|_\op \leq \epsilon \;\;\Longrightarrow\;\; (1-\epsilon)(T+\lambda) \preceq T_n+\lambda \preceq (1+\epsilon)(T+\lambda)
$$
and subsequently $\|(T+\lambda)^{1/2}(T_n+\lambda)^{-1/2}\|_\op \leq 1/\sqrt{1-\epsilon}$.
}
This proves the first property, because Eq.~(\ref{eq:proof-gdiff-0}) is uniform almost everywhere in $z$. This proves the first property, because Eq.~(\ref{eq:proof-gdiff-0}) is uniform in $z$.

% \pd{The upper bound on $\|(T+\lambda)^{-1/2}(T_n-T)(T_n+\lambda)^{-1/2}\|_\op$ does not follow from Lemma 8.1 directly, as Lemma 8.1 provides a bound for the symmetric whitening $\|(T+\lambda)^{-1/2}(T_n-T)(T+\lambda)^{-1/2}\|_\op$ but this step needs mixed whitening with $(T_n+\lambda)^{-1/2}$. Perhaps a negligible term will creep into the exponent. We need to add details on how this term is controlled.}

For the second identity, because $\{g(z,Z_i)^2\}_{i=1}^n$ are i.i.d.~random variables,
the Chebyshev's inequality asserts that for any $t>0$, 
\begin{align}
\Pr\left[\left|\frac{1}{n}\sum_{i=1}^n \bar g(z,Z_i)^2 - \mathbb E[g(z,Z)^2]\right| > t\mathbb E[g(z,Z)^2]\right] \leq \frac{\mathbb E[g(z,Z)^4]}{nt^2(\mathbb E[g(z,Z)^2])^2} \leq \frac{C_p}{nt^2},
\label{eq:proof-gdiff-2}
\end{align}
where the last inequality holds thanks to Assumption (B2). 
{\color{magenta}
Because the right-hand side of Eq.~(\ref{eq:proof-gdiff-2}) is uniform in almost every $z$,
This implies that for almost every $z$,
\begin{equation*}
\frac{1}{n}\sum_{i=1}^n\bar g(z,Z_i)^2 \;\;\overset{p}{\to}\;\; \mathbb E[g(z,Z)^2],
\end{equation*} 
This completes the proof. 
}

{\color{magenta} YW: check if this statement is accurate?}
\end{proof}

\begin{lemma}
Assume the same set of conditions in Lemmas \ref{lem:empirical-population-cov} and \ref{lem:gdiff}. 
Let $\varepsilon_i = Y_i - \langle X_i,\beta_0\rangle$. 
Then for any $x\in\mX$ and $z=L_{K^{1/2}}x$, as $n\to\infty$, 
$$
\frac{1}{\tilde\sigma_n(z)\sqrt{n}}\sum_{i=1}^n \varepsilon_i g(z,Z_i)\;\;\overset{d}{\to}\;\; \mathcal N(0,1),
$$
where $Z_i=L_{K^{1/2}}X_i$, $g(z,Z)=\langle z,(T_n+\lambda_n I)^{-1}Z\rangle$ and $\tilde\sigma_n(z)^2 = n^{-1}\sum_{i=1}^n \sigma_\varepsilon^2g(z,Z_i)^2$.
\label{lem:isoclt}
\end{lemma}

\begin{proof}[Proof of Lemma \ref{lem:isoclt}.]
Define $\bar g(z,Z) := \langle z,(T+\lambda_n I)^{-1}Z\rangle$.
Because $\{\varepsilon_i\bar g(z,Z_i)\}_{i=1}^n$ are i.i.d.~centered random variables with finite variance, the standard central-limit theorem yields that
\begin{equation}
\frac{1}{\bar\sigma(z)\sqrt{n}}\sum_{i=1}^n \varepsilon_i\bar g(z,Z_i) \;\;\overset{d}{\to}\;\;\mathcal N(0,1),
\label{eq:proof-isoclt-1}
\end{equation}
where
$$
\bar\sigma(z) =\sigma_\varepsilon^2\mathbb E[g(z,Z)^2]= \sigma_\varepsilon^2\mathbb E[\langle z, (T+\lambda_n I)^{-1}Z\rangle^2]
$$

{\color{cyan}
Let $Z,Z'$ be arbitrary in the support of $P_{\mathcal Z}$ induced by $P_{\mathcal X}$ with the understanding that $Z=L_{K^{1/2}}X$,
and for every $j\geq 1$ define $\alpha_j := \langle Z,\psi_j\rangle$, $\alpha_j':=\langle Z',\psi_j\rangle$,
where $\{\psi_j\}_{j=1}^{\infty}$ is the orthonormal basis of the population covariance operator $T=\mathbb E[Z\otimes Z]$.
Thanks to Assumption (B3), it holds almost surely that $|\alpha_j|,|\alpha_j'|\leq C_\nu n^{\kappa}\sqrt{s_j}$ for every $j\geq 1$.
Subsequently, with probability 1
$$
\big|\langle Z,(T+\lambda_n)^{-1}Z\rangle\big| \leq C_\nu n^{2\kappa}\sum_{j=1}^{\infty} \frac{s_j}{s_j+\lambda_n}\lesssim n^{2\kappa}\times \int_1^{\infty}\frac{ j^{-2r}\ud j}{j^{-2r}+\lambda_n} \lesssim n^{2\kappa+1/(2r)+\delta}
$$
for any $\delta>0$, where the last inequality holds because $\lambda_n n\to 0$ and $\lambda_n n^{1+\delta}\to \infty$ for any $\delta>0$.
}

\pd{This is a triangular array CLT because of $\lambda_n$, hence Lyapunov/Lindeberg verification is needed.}
{\color{magenta}YW: this isn't easy. We need, for almost every $z$,
$$
\lim_{n\to\infty}\frac{1}{\sqrt{n}}\frac{\mathbb E[|\langle z, (T+\lambda_n)^{-1}Z\rangle|^3]}{\mathbb E[\langle z,(T+\lambda_n)^{-1}Z\rangle^2]^{3/2}} = 0.
$$
This is not straightforward and would depend on the actual structure of $Z$. The smoothing effects from $\lambda_n$ would not be sufficient for it to hold unconditionally. Perhaps we add an assumption about this?

If we use Linderberg we need similar assumptions, controlling behaviors when $\langle z,(T+\lambda_n)^{-1} Z\rangle$ is very large.
}

Using Cauchy-Schwarz inequality, it holds almost surely that
\begin{align*}
\left|\frac{1}{n}\sum_{i=1}^n g(z,Z_i)^2-\bar g(z,Z_i)^2\right| &\leq \left|\frac{2}{n}\sum_{i=1}^n [g(z,Z_i)-\bar g(z,Z_i)]\bar g(z,Z_i)\right|\\
&\leq 2\sqrt{\frac{1}{n}\sum_{i=1}^n[g(z,Z_i)-\bar g(z,Z_i)]^2}\sqrt{\frac{1}{n}\sum_{i=1}\bar g(z,Z_i)^2}.
\end{align*}
Lemma \ref{lem:gdiff} then asserts that the first term on the right-hand side of the above inequality converges in probability to zero, and the second term converges in probability to $\sqrt{\mathbb E[g(z,Z_i)^2]}$.
Therefore,
\begin{equation}
\tilde\sigma_n(z)^2 \;\;\overset{p}{\to}\;\; \bar\sigma(z)^2.
\label{eq:proof-isoclt-2}
\end{equation}

Next, using Chebyshev's inequality, for any $\delta>0$, conditioned on $\{Z_i\}$ it holds with probability $1-\delta$ that
$$
\left|\frac{1}{\sqrt{n}}\sum_{i=1}^n\varepsilon_i[g(z,Z_i)-\bar g(z,Z_i)]\right| \leq \frac{\sigma_\varepsilon^2}{\delta^2}\sqrt{\frac{1}{n}\sum_{i=1}^n[g(z,Z_i)-\bar g(z,Z_i)]^2} 
$$
Because $\frac{1}{n}\sum_{i=1}^n[g(z,Z_i)-\bar g(z,Z_i)]^2\overset{p}{\to}0$, for any $\epsilon,u>0$ there exists $N\in\mathbb N$ such that for any $n\geq N$, $\frac{1}{n}\sum_{i=1}^n[g(z,Z_i)-\bar g(z,Z_i)]^2\leq\epsilon$
with probability $1-u$.
Now for any $\epsilon',\delta'>0$, pick $\delta=\delta'/2$, $u=\delta'/2$ and $\epsilon=(\epsilon')^2\delta^4/\sigma_\varepsilon^4>0$, there exists $N\in\mathbb N$ such that for any $n\geq N$, 
$$
\Pr\left[\left|\frac{1}{\sqrt{n}}\sum_{i=1}^n\varepsilon_i[g(z,Z_i)-\bar g(z,Z_i)]\right|>\epsilon'\right]\leq\delta',
$$
where the probability is on both $\{\varepsilon_i\}$ and $\{Z_i\}$, and the above inequality holds by the union bound.
This implies that
\begin{equation}
\frac{1}{\bar\sigma(z)\sqrt{n}}\sum_{i=1}^n\varepsilon_i[g(z,Z_i)-\bar g(z,Z_i)] \;\;\overset{p}{\to}\;\; 0.
\label{eq:proof-isoclt-3}
\end{equation}
Combining Eqs.~(\ref{eq:proof-isoclt-1},\ref{eq:proof-isoclt-2},\ref{eq:proof-isoclt-3}), the proof of Lemma \ref{lem:isoclt} is complete with the Slutsky's theorem.
\end{proof}


\subsection{Proof of Theorem \ref{thm:isokrr-limiting}}

To simplify notations, in the proofs of Theorems \ref{thm:isokrr-limiting} and \ref{thm:isokrr-rates} we shall drop all $\iso$ superscripts as this is the only set of estimators, confidence intervals and their related
quantities to be analyzed in the proofs.
With $Y_i=\langle X_i,\beta_0\rangle+\varepsilon_i = \langle Z_i,\theta_0\rangle+\varepsilon_i$, it holds that
$$
\hat\theta_n = (T_n+\lambda_n)^{-1}T_n\theta_0 + \frac{1}{n}\sum_{i=1}^n\varepsilon_i (T_n+\lambda_n)^{-1}Z_i.
$$
Subsequently,
\begin{align}
\hat\theta_n-\theta_0 &= \underbrace{-\lambda_n (T_n+\lambda_n)^{-1}\theta_0}_{=: b_n} + \underbrace{\frac{1}{n}\sum_{i=1}^n\varepsilon_i(T_n+\lambda_n)^{-1}Z_i}_{=:v_n}.
\end{align}

For any test function $x\in\supp(P_X)$, let $z=L_{K^{1/2}}x=\sum_{j=1}^{n}\hat\nu_j\hat\psi_j$, where $\hat\nu_j=\langle z,\hat\psi_j\rangle$.
Let also $\hat\vartheta_j = \langle\theta_0,\hat\psi_j\rangle$.
Slightly abusing notations by denoting $b_n(z) := \langle b_n,z\rangle = \int_I b_n(t)z(t)\ud t$, we have that
\begin{align}
\big|b_n(z)\big| &= \lambda_n\left|\left\langle(T_n+\lambda_n)^{-1}\sum_{j=1}^n\hat\vartheta_j\hat\psi_j, \sum_{j=1}^n\hat\nu_j\hat\psi_j\right\rangle\right|
= \lambda_n \left|\sum_{j=1}^n \frac{\hat\nu_j\hat\vartheta_j}{\hat s_j+\lambda_n}\right|\label{eq:proof-simple-1}\\
&\leq \lambda_n \sqrt{\sum_{j=1}^n\hat\vartheta_j^2}\sqrt{\sum_{j=1}^n\frac{\hat\nu_j^2}{(\hat s_j+\lambda_n)^2}} \label{eq:proof-simple-2}\\
&\leq \sqrt{\lambda_n}\sqrt{\sum_{j=1}^n\hat\vartheta_j^2}\sqrt{\sum_{j=1}^n\frac{\hat\nu_j^2}{\hat s_j+\lambda_n}}\label{eq:proof-simple-2half}\\
&= \sqrt{\lambda_n}\cdot \sqrt{\langle z,(T_n+\lambda_n)^{-1}z\rangle} = {\frac{\sqrt{\lambda_nn}}{\sigma_\varepsilon}\hat\sigma_n(z)}.\label{eq:proof-simple-3}
\end{align}
In the above derivations, Eq.~(\ref{eq:proof-simple-1}) holds because $\{\hat\psi_j\}_{j=1}^n$ is orthonormal with respect to $L_2(I)$ and $T_n=\sum_{j=1}^n\hat s_j\hat\psi_j\otimes\hat\psi_j$;
Eq.~(\ref{eq:proof-simple-2}) holds by the Cauchy-Schwarz inequality;
Eq.~(\ref{eq:proof-simple-2half}) holds by noting that $\lambda_n/(\sqrt{\hat s_j+\lambda_n})\leq\sqrt{\lambda_n}$ for all $j$, and furthermore
$$
\langle z, (T_n+\lambda_n)^{-1}z\rangle = \left\langle\sum_{j=1}^n\hat\nu_j\hat\psi_j, \sum_{j=1}^n\frac{\hat\nu_j}{\hat s_j+\lambda_n}\hat\psi_j\right\rangle = \sum_{j=1}^n\frac{\hat\nu_j^2}{\hat s_j+\lambda_n}.
$$
With $\lambda_n n\to 0$ as $n\to\infty$, Eq.~(\ref{eq:proof-simple-3}) implies that
$$
b_n(z)/\hat\sigma_n(z)\;\;\to\;\;0\;\;\;\;\;\;a.s.
$$

We next analyze the term involving $v_n$ on an arbitrary test function $x\in\supp(P_X)$ with $z=L_{K^{1/2}}x=\sum_{j=1}^n\hat\nu_j\hat\psi_j$, where $\hat\nu_j=\langle z,\hat\psi_j\rangle$.
Let $v_n(z)=\langle v_n,z\rangle = \int_I v_n(t)z(t)\ud t$. We have that
$$
v_n(z) = \frac{1}{n}\sum_{i=1}^n \varepsilon_i \langle z, (T_n+\lambda_n)^{-1}Z_i\rangle.
$$
Note that $v_n(z)$ is \emph{pivotal}, meaning that it does not depend on the unknown $\beta_0\in\Theta_2(\mH)$ or $\theta_0=L_{K^{1/2}}\beta_0$.
Note also that $\{\varepsilon_i\}_{i=1}^n$ are conditionally centered when $\{Z_i\}_{i=1}^n$ are fixed. Therefore, if $\{\varepsilon_i\}_{i=1}^n$ are Normally distributed,
the distribution of $v_n(z)$ is an exact Normal distribution with mean zero and variance equal to
$$
\frac{{\sigma_\varepsilon^2}}{n^2}\sum_{i=1}^n \langle z,(T_n+\lambda_n)^{-1}Z_i\rangle^2 = \frac{{\sigma_\varepsilon^2}}{n}\langle z, (T_n+\lambda_n)^{-1}T_n(T_n+\lambda_n)^{-1}z\rangle = \frac{\tilde\sigma_n(z)^2}{n}
$$
because $T_n=\frac{1}{n}\sum_{i=1}^n Z_i\otimes Z_i$. If $\{\varepsilon_i\}_{i=1}^n$ are not Normally distributed, Lemma \ref{lem:isoclt} asserts that under imposed conditions,
$$
\frac{1}{\tilde\sigma_n(z)\sqrt{n}}\sum_{i=1}^n\varepsilon_i\langle z,(T_n+\lambda_n)^{-1}Z_i\rangle\;\;\overset{d}{\to}\;\; \mathcal N(0,1),
$$
which is to be proved.

\subsection{Proof of Theorem \ref{thm:isokrr-rates}}

For any $x\in\supp(P_X)$ and $z=L_{K^{1/2}}x$, the definition of $\varphi(x)$ yields
%$$
%n\big|\varphi(x)\big|^2 = \sigma_\varepsilon^2\langle z, (T_n+\lambda_n)^{-1}z\rangle^2.
%$$
$$
n\big|\varphi(x)\big|^2 = 4 \tau^2_{\alpha/2}\sigma_\varepsilon^2\langle z, (T_n+\lambda_n)^{-1}z\rangle.
$$
Because $|\varphi(x)|$ is pivotal, we may safely drop the $\sup_{\beta_0\in\Theta_2(\mH)}$ notation and instead calculate
\begin{align*}
n\mathbb E_{P_X}\left[|\varphi(X)|^2\right] &= n\mathbb E_{P_Z}[\langle Z,(T_n+\lambda_n)^{-1}Z\rangle^2] = \sigma_\varepsilon^2\tr((T_n+\lambda_n)^{-1/2}T(T_n+\lambda_n^{-1/2}))\\
&= \sigma_\varepsilon^2(1+o_P(1))\tr((T+\lambda_n)^{-1/2}T(T+\lambda_n)^{-1/2}),
\end{align*}
where the last identity holds by invoking Lemma \ref{lem:empirical-population-cov}. Additionally, invoking Assumption (A1) and the fact that $\lambda_n\gtrsim n^{-1-\delta}$ for any $\delta>0$, it holds that
\begin{align*}
\tr((T+\lambda_n)^{-1/2}T(T+\lambda_n^{-1/2})) &= \sum_{j=1}^{\infty} \frac{s_j}{s_j+\lambda_n} \lesssim \int_1^{\infty}\frac{j^{-2r}}{j^{-2r}+\lambda_n}\ud j\lesssim n^{1/2r+\delta},
\end{align*}
which is to be proved.

%\pd{Just realized that $I$ is missing in $T_n+\lambda_n I$ in many places. Needs to be fixed.}

%\pd{I did not check this rigorously, but we should make sure that conditioning on high-probability spectral events is properly accounted for in the final honesty statements. Overall, the details look reasonable, but this point needs to be verified carefully for rigor.}

\subsection{Proof of Theorem \ref{thm:isokrr-lb}}

In this lower bound we work with $\gamma_0=0$ and the definition of honest confidence intervals that
$$
\Pr_{\beta_0}[\langle x,\beta_0\rangle\in\varphi(x)] \geq 1-\alpha,\;\;\;\;\;\;\forall \beta_0\in\Theta_2(\mH),\;\;x\in\supp(P_X).
$$
The definition in Eq.~(\ref{eq:honest-ci}) is asymptotic in nature, and would imply the above property for a slightly smaller constant significance level for sufficiently large $n$.

Let $\{\psi_j\}_{j=1}^{\infty}$ be an arbitrary orthonormal basis of $L_2(I)$.
Consider a distribution $P_X$ such that for every $X\sim P_X$, $\langle X,\psi_j\rangle$ are \emph{independently distributed} random variables that satisfy
$$
\Pr[\langle X,\psi_j\rangle = \sqrt{s_j}] = \Pr[\langle X,\psi_j\rangle = -\sqrt{s_j}] = \frac{1}{2},\;\;\;\;\;\;\forall j\geq 1.
$$
with
$$
s_j = \frac{1}{2r-1}\times j^{-2r}.
$$
Let $\{\varepsilon_i\}_{i=1}^n$ be i.i.d.~standard Normally distributed random variables. It is easy to verify that this problem instance satisfies Assumptions (A1), (A3), (B1), (B2), (B3),
with $\sigma_\varepsilon=1$, $M_\varepsilon=1.6$, $c_s=1/(2r-1)$, $C_s=1$, $C_p=2$, $\kappa=0$ and $C_\nu=1$.

Let $M$ be an integer that grows in a particular way with $n$, to be specified later in this proof. 
Consider a model class $\Theta_n$ such that for every $\beta_0\in\Theta_n$, either $\beta_0\equiv 0$ or
$$
\langle L_{K^{-1/2}}\beta_0, \psi_j\rangle = \left\{\begin{array}{ll} \pm 1/\sqrt{M},& \text{if }M< j\leq 2M;\\ 0,& \text{otherwise}.\end{array}\right.
$$
It is easy to see that $\Theta_n\subseteq\Theta_2(\mH)$. Furthermore, all $\beta_0\in\Theta_n$ would satisfy Assumption (A2).
There are a total of $2^M+1$ hypothetical regression functionals (the one additional functional comes from $\beta_0\equiv 0$) in $\Theta_n$.
For convenience, given any binary vector $\vct\iota\in\{\pm 1\}^{2M}$, define
$$
\beta_{\vct\iota} := \sum_{j=M+1}^{2M}\frac{\vct\iota_j}{\sqrt{M}} L_{K^{1/2}}\psi_j.
$$
Additionally, we use $\beta_o\equiv 0$ to denote the zero regression functional.

Because $\varphi$ is an honest confidence interval at significance level $\alpha$, it holds that 
\begin{equation}
\Pr_{\beta_o}\left[0\in\varphi(x)\right] \geq 1-\alpha,\;\;\;\;\;\;\forall x\in\supp(P_X).
\label{eq:proof-isolb-1}
\end{equation}

Fix arbitrary $\vct\iota\in\{\pm 1\}^{2M}$, let $x_{\vct\iota}\in\supp(P_X)$ be such that 
\begin{align}
\sgn(\langle x_{\vct\iota},\psi_j\rangle)=\vct\iota_j,\;\;\;\;\;\;\forall j\leq 2M\label{eq:proof-isolb-xl}
\end{align}
and that $\langle x_{\vct\iota},\psi_j\rangle = 0$ for all $j>2M$.
Note that there might be multiple such test functions satisfying the sign matching property: an arbitrary test function could be selected.
Define constant $a_n$ as
$$
a_n := \langle x_{\vct\iota},\beta_{\vct\iota}\rangle = \frac{1}{2r-1}\sum_{j=M+1}^{2M}\frac{j^{-r}}{\sqrt{M}}
$$
which satisfies
\begin{equation}
\frac{(2M)^{-r}\sqrt{M}}{2r-1} \leq a_n \leq \frac{M^{-r}\sqrt{M}}{2r-1}.
\label{eq:proof-isolb-1half}
\end{equation}
Then, invoking again the honesty of $\varphi$, it holds that
\begin{equation}
\Pr_{\beta_{\vct\iota}}\left[a_n\in\varphi(x_{\vct\iota})\right] \geq 1-\alpha.
\label{eq:proof-isolb-2}
\end{equation}
Furthermore, for any $\beta_{\vct\iota}\in\Theta_2(\mH)$, 
\begin{align}
\mathrm{KL}(P_{\beta_o}^n\|P_{\beta_{\vct\iota}}^n) &= n\times \mathbb E_{X\sim P_X}\mathbb E_{\beta_o}\left[\log\frac{P(Y|X,\beta_{\vct\iota})}{P(Y|X,\beta_o)}\right]
= \frac{n}{2}\mathbb E_{X\sim P_X}[\langle X, \beta_{\vct\iota}-\beta_o\rangle^2]\nonumber\\
&= \frac{n}{2}\sum_{j=M+1}^{2M} \frac{1}{M}\times \frac{j^{-2r}}{(2r-1)^2} \leq \frac{nM^{-2r}}{2(2r-1)^2}.\label{eq:proof-isolb-3}
\end{align}
Selecting 
$$
M\asymp n^{1/2r}
$$
the right-hand side of Eq.~(\ref{eq:proof-isolb-3}) is upper bounded by 0.1, and therefore using Pinsker's inequality,
\begin{align}
{\Pr_{\beta_o}[a_n\in\varphi(x_{\vct\iota})]\geq \Pr_{\beta_{\vct\iota}}[a_n\in\varphi(x_{\vct\iota})] - \sqrt{2\mathrm{KL}(P_{\beta_o}^n\|P_{\beta_{\vct\iota}}^n)} \geq 1/2.}
\label{eq:proof-isolb-4}
\end{align}
{\color{magenta}
Because Eq.~(\ref{eq:proof-isolb-4}) holds for any $\vct\iota\in\{\pm 1\}^{2M}$, in conjunction with Eq.~\eqref{eq:proof-isolb-1} we have that for each $\vct\iota\in\{\pm 1\}^{2M}$ and the context vector $x_{\vct\iota}$ 
defined in Eq.~(\ref{eq:proof-isolb-xl}) it holds that
\begin{align*}
    E_{\beta_o}[|\varphi(x_{\vct\iota})|^2] \geq a_n^2/2.
\end{align*}
Because $\{x_{\vct\iota}\}_{\vct\iota\in\{\pm 1\}^2M}$ uniformly explores the support of $P_X$ that is also by definition a uniform distribution, it holds that
\begin{align*}
{\mathbb E_{X\sim P_X} \mathbb E_{\beta_o}[|\varphi(X)|^2] } &\geq a_n^2/2 \gtrsim M^{1-2r} \asymp n^{-(2r-1)/2r},
\end{align*}
}
which is to be proved.

%\pd{We need to reason out why we have the lower bound in expectation from the fact that for each $\vct\iota\in\{\pm 1\}^{2M}$ there exists $x_{\vct\iota}$ such that $ E_{\beta_o}[|\varphi(x_{\vct\iota})|^2] \geq a_n^2/2$. This is because the lower bound need not hold for all support points $x_{\vct\iota}$ but rather some support points where it is true, and it is not clear whether $P_X$ has enough mass on those support points. One option is to say that the expectation is lower bounded $E_{\beta_o}[|\varphi(x_{\vct\iota})|^2]/(\# \text{of support points})$ for some $\vct\iota$. Does this work?}


\subsection{Proof of Theorem \ref{thm:anikrr-limiting}}

In the proofs of Theorems \ref{thm:anikrr-limiting} we shall drop all $\ani$ superscripts as this is the only set of estimators, confidence intervals and their related quantities to be analyzed in these proofs.
Following similar derivations that lead to Eq.~(\ref{eq:proof-isoclt-1}), and accounting for the null space of $\{\hat\psi_j\}_{j=1}^n$, we have for the anisotropic case that
\begin{equation}
\hat\theta_n-\theta_0 =\underbrace{-(I-\hat\Pi)\theta_0}_{=:t_n}\underbrace{- (T_n+\Lambda_n)^{\dagger}\Lambda_n\theta_0}_{=: b_n} + \underbrace{\frac{1}{n}\sum_{i=1}^n\varepsilon_i(T_n+\Lambda_n)^{\dagger}Z_i}_{=:v_n},
\label{eq:proof-aniclt-1}
\end{equation}
where $\hat\Pi = \sum_{j=1}^n\hat\psi_j\otimes\hat\psi_j$.

Consider $X_1$ being an arbitrary input function in $\{(X_i,Y_i)\}_{i=1}^n$,
and let $Z_1=L_{K^{1/2}}X_1$, $b_n(Z_1)=\langle b_n,Z_1\rangle$, $v_n(Z_1)=\langle v_n,Z_1\rangle$.
Noet that $t_n(Z_1)=0$ because $\hat\Pi(Z_1)=Z_1$.
Because $v_n(z)$ is either exactly Normally distributed (if $\{\varepsilon_i\}_{i=1}^n$ are Normally distributed) \pd{conditional on $Z_i$}, or satisfies Lyapnov conditions of central limit theorem with probabilities converging to one
provided that Eq.~(\ref{eq:anikrr-lyapnov}) is satisfied, we know that $v_n(z)$ is asymptotically Normally distributed whose variance is exactly $\frac{\sigma_\varepsilon^2}{n^2}\sum_{i=1}^n\langle z, (T_n+\Lambda_n)^{\dagger}Z_i\rangle^2 = \frac{\sigma_\varepsilon^2}{n}\langle z, (T_n+\Lambda_n)^{\dagger}T_n(T_n+\Lambda_n)^{\dagger}z\rangle=\hat\sigma_n^\ani(z)^2$. This shows
\begin{equation}
\frac{v_n(z)}{\hat\sigma_n^\ani(z)}\;\;\overset{d}{\to}\;\;\mathcal N(0,1).
\label{eq:proof-aniclt-2}
\end{equation}

We next analyze the bias term. Using the spectral decomposition of both $T_n$ and $\Lambda_n$, $b_n(z)$ can be expressed as
$$
b_n(z) = \left\langle(T_n+\Lambda_n)^{\dagger}\Lambda_n\sum_{j=1}^n\hat\vartheta_j\hat\psi_j, \sum_{j=1}^n\hat\nu_j\hat\psi_j\right\rangle = \sum_{j=1}^n\frac{\lambda_{nj}\hat\nu_j\hat\vartheta_j}{\hat s_j+\lambda_{nj}},
$$
where $\hat\vartheta_j=\langle\theta_0,\hat\psi_j\rangle$ and $\hat\nu_j = \langle z,\hat\psi_j\rangle$.
Using Cauchy-Schwarz inequality and the fact that $\sum_{j=1}^n\hat\vartheta_j^2\leq 1$ because $\|\theta_0\|_2^2 = \int_I\theta_0(t)^2\ud t\leq 1$, it holds that
$$
b_n(z)^2 \leq \sum_{j=1}^n\frac{\lambda_{nj}^2\hat\nu_j^2}{(\hat s_j+\lambda_{nj})^2}.
$$
Using the choice that $\lambda_{nj}= \sqrt{\hat s_j\lambda_n}$, the above inequality simplifies to
\begin{align}
b_n(z)^2 \leq \lambda_n\sum_{j=1}^n\frac{\hat\nu_j^2(\lambda_n\wedge\hat s_j)}{(\hat s_j+(\lambda_n\wedge\sqrt{\hat s_j\lambda_n}))^2}\leq \lambda_n\sum_{j=1}^n\frac{\hat\nu_j^2\hat s_j}{(\hat s_j+\sqrt{\hat s_j\lambda_n})^2}.
\label{eq:proof-aniclt-3}
\end{align}
On the other hand, expanding $\hat\sigma_n^2(z)$ in the spectrum of $\{\hat\psi_j\}_{j=1}^n$ we obtain
\begin{align}
\hat\sigma_n^2(z) &= \frac{\sigma_\varepsilon^2}{n}\langle z, (T_n+\Lambda_n)^{\dagger}T_n(T_n+\Lambda_n)^{\dagger}z\rangle = \frac{\sigma_\varepsilon^2}{n}\sum_{j=1}^n\frac{\hat\nu_j^2\hat s_j}{(\hat s_j+\sqrt{\hat s_j\lambda_n})^2}.
\label{eq:proof-aniclt-4}
\end{align}
Comparing Eqs.~(\ref{eq:proof-aniclt-3},\ref{eq:proof-aniclt-4}) and noting that $\lambda_n n\to0$ as $n\to\infty$, we obtain
\begin{equation}
\frac{|b_n(z)|}{\hat\sigma_n(z)} \leq \sigma_\varepsilon^{-1}\sqrt{\lambda_n n} \to 0\;\;\;\;\;a.s.
\label{eq:proof-aniclt-5}
\end{equation}
Eqs.~(\ref{eq:proof-aniclt-1},\ref{eq:proof-aniclt-5}) together yield Theorem \ref{thm:anikrr-limiting}, by Slutsky's theorem.

\subsection{Proof of Theorem \ref{thm:anikrr-rates}}

Let $Z_i=L_{K^{1/2}}X_i$.
It holds that
\begin{align*}
\frac{1}{n}\sum_{i=1}^n\big|\varphi(X_i)\big|^2 &= \frac{4\tau_\alpha^2}{n^2}\sum_{i=1}^n\hat\sigma_n^\ani(z_i)^2 = \frac{4\sigma_\varepsilon^2\tau_\alpha^2}{n^2}\sum_{i=1}^n \langle Z_i, (T_n+\Lambda_n)^{\dagger}T_n(T_n+\Lambda_n)^{\dagger}Z_i\rangle\\
&= \frac{4\sigma_\varepsilon^2\tau_\alpha^2}{n}\tr(((T_n+\Lambda_n)^{\dagger}T_n)^2) = \frac{4\sigma_\varepsilon^2\tau_\alpha^2}{n}\sum_{j=1}^n\frac{\hat s_j^2}{(\hat s_j+\lambda_{nj})^2}\\
& \overset{(a)}{\leq} \frac{4\sigma_\varepsilon^2\tau_\alpha^2}{n}\sum_{j=1}^n\frac{\hat s_j}{(\sqrt{\hat s_j}+\sqrt{\lambda_n})^2} \leq \frac{4\sigma_\varepsilon^2\tau_\alpha^2}{n}\sum_{j=1}^n\frac{\hat s_j}{{\hat s_j}+{\lambda_n}} \\
&\lesssim \frac{1}{n}\int_1^{\infty} \frac{j^{-2r}}{j^{-2r}+\lambda_n}\ud j \lesssim n^{-\frac{2r-1}{2r}+\delta}
\end{align*}
for any $\delta>0$, where Eq.~(a) holds because $\lambda_{nj}= \sqrt{\hat s_j\lambda_n}$, and
the last inequality holds by the condition imposed on $\hat\{s_j\}$ and the asymptotic scaling that $\lambda_n\gtrsim n^{-1-\delta}$ for any $\delta>0$.

{\color{red} YW: please check again. this derivation should be correct and no longer need to do $\lambda_n\wedge\sqrt{\hat s_j\lambda_n}$.}




\subsection{Proof of Lemma \ref{lem:stability}}

Express both $T$ and $T_n$ on the eigenfunctions $\{\psi_j\}_{j=1}^n$ and obtain their infinite-dimensional matrix representations $A$ and $A_n$,
where $A_{ij}=\langle \psi_i,T\psi_j\rangle$ and $[A_n]_{ij}=\langle\psi_i, T_n\psi_j\rangle$ for all $i,j\geq 1$.
More specifically,
$$
A_{ij} = \mathbb E[\langle Z,\psi_i\rangle\langle Z,\psi_j\rangle] = \vct 1\{i=j\}s_i.
$$
On the other hand, $[A_n]_{ij}$ is the sample average of a sum of $n$ i.i.d.~random variables and is potentially non-zero for $i\neq j$:
$$
[A_n]_{ij} = \frac{1}{n}\sum_{\ell=1}^n \langle Z_\ell,\psi_i\rangle\langle Z_\ell,\psi_j\rangle.
$$

%{\color{blue}
%Consider the submatrix of $A_n$ corresponding to eigenfunctions $J(n)+1,\cdots,n$. The trace formula and Eq.~(\ref{eq:proof-stability-1}) implies that, with probability $1-n^{-1}$,
%\begin{align}
%\sum_{j=J(n)+1}^n\hat s_j &= \sum_{j=J(n)+1}^n [A_n]_{jj} \leq \sum_{j=J(n)+1} s_j + O\left(\frac{s_j\sqrt{\log n}}{\sqrt{n}} + \frac{s_jn^{2\kappa}\log n}{n}\right)\nonumber\\
%&\lesssim \int_{J(n)}^{\infty} j^{-2r}\left(1+\frac{\sqrt{\log n}}{\sqrt{n}}\right)\ud j  \lesssim J(n)^{-2r+1}.
%\end{align}
%This proves the second property of Lemma \ref{lem:stability}.
%}

%{\color{blue}
%Consider the submatrix of $A_n$ corresponding to eigenfunctions $J(n)+1,\cdots,n$. The trace formula yields
%\begin{align}
%\sum_{j=J(n)+1}^n\hat s_j-s_j &= \sum_{j=J(n)+1}^{\infty} [A_n]_{jj}-A_{jj} = \sum_{j=J(n)+1}^n \langle\psi_j, (T_n-T)\psi_j\rangle.
%\label{eq:proof-stability-1-1}
%\end{align}
%For every $i\in[n]$, define $w_i := \sum_{j=J(n)+1}^{\infty}\langle\psi_j, Z_i\rangle^2$ such that $\mathbb E[w_i] = \sum_{j=J(n)+1}^{\infty}s_j$.
%The right-hand side of Eq.~(\ref{eq:proof-stability-1-1}) is precisely $\frac{1}{n}\sum_{i=1}^n w_i-\mathbb E[w_i]$.
%Furthermore, 
%\begin{align*}
%\mathbb E[w_i^2] &=\sum_{j,k>J(n)}\mathbb E[\langle Z_i,\psi_j\rangle^2\langle Z_i,\psi_k\rangle^2]\overset{(a)}{\leq} \left(\sum_{j>J(n)}\sqrt{\mathbb E[\langle Z_i,\psi_j\rangle^4]}\right)^2\\
%&\overset{(b)}{\leq} \left(\sum_{j>J(n)}C_ps_j\right)^2 \overset{(c)}{\lesssim} \left(\int_{J(n)}^{\infty} j^{-2r}\ud j\right)^2 \lesssim J(n)^{2-4r},
%\end{align*}
%where Eq.~(a) holds by the Cauchy-Schwarz inequality, Eq.~(b) holds by Assumption (B2) and Eq.~(c) holds by Assumption (B1).
%Additionally, Assumption (A3) implies that $w_i\leq \|Z\|_2^2\leq 1$ almost surely. Subsequently, using Bernstein's inequality, with probability $1-n^{-1}$ it holds that
%\begin{align}
%\left|\sum_{j=J(n)+1}^n \langle\psi_j, (T_n-T)\psi_j\rangle\right| &= \left|\frac{1}{n}\sum_{i=1}^n w_i-\mathbb E[w_i]\right| \lesssim \frac{J(n)^{1-2r}\sqrt{\log n}}{\sqrt{n}} + \frac{\log n}{n}.
%\label{eq:proof-stability-1-2}
%\end{align}
%Combining Eqs.~(\ref{eq:proof-stability-1-1},\ref{eq:proof-stability-1-2}) and noting that $\sum_{j>J(n)}s_j \lesssim \int_{J(n)}^{\infty} j^{-2r}\ud j \lesssim J(n)^{1-2r}$, we complete the proof of the second property of Lemma \ref{lem:stability}.
%}


Now, let $C=(2C_s/c_s)^{1/(2r)}>0$ be a constant such that, for every $j\in\mathbb N$, $s_{Cj}\leq C_s (Cj)^{-2r} \leq c_sj^{-2r}/2\leq s_j/2$ thanks to Assumption (B1). 
Let $\bar J=\lceil CJ(n)\rceil$.
By Assumption (B2) and the Cauchy-Schwarz inequality, it holds for every $i,j\geq 1$ that 
$$
\mathbb E[\langle Z,\psi_i\rangle^2\langle Z,\psi_j\rangle^2] \leq \sqrt{\mathbb E[\langle Z,\psi_i\rangle^4]}\sqrt{\mathbb E[\langle Z,\psi_j\rangle^4]}
\leq C_p\mathbb E[\langle Z,\psi_i\rangle^2]\mathbb E[\langle Z,\psi_j\rangle^2] = C_ps_is_j.
$$
By Assumption (B3), it holds for every $i,j\geq 1$ that 
$$
|\langle Z,\psi_i\rangle\langle Z,\psi_j\rangle| \leq C_\nu^2\sqrt{s_js_k}n^{2\kappa}
$$ almost surely.
Subsequently, using Bernstein's inequality, for every $i,j\geq 1$, with probability $1-n^{-c}/(3i^2j^2)$ it holds that 
\begin{align}
\big|[A_n]_{ij}-A_{ij}\big| \leq C_1\left(\sqrt{\frac{C_ps_is_j\log (nij)}{n}} + \frac{C_\nu\sqrt{s_is_j}n^{2\kappa}\log (nij)}{n}\right) \lesssim \frac{\sqrt{s_is_j}\log(nij)}{\sqrt{n}},
\label{eq:proof-stability-1}
\end{align}
where $C_1<+\infty$ is a numerical constant depending only on $c$, and the last inequality holds provided that $\kappa<1/4$.
Using the union bound (i.e.~Bonferroni's inequality) and noting that $(\sum_{j=1}^{\infty}1/j^2)^2=\pi^4/36<3$, with probability $1-n^{-c}$ Eq.~(\ref{eq:proof-stability-1}) holds for all $i,j\geq 1$.

Let $e_1,e_2,\cdots$ be the standard basis vectors.
We will prove that, for every $j\in[\bar J]$, there exists $w$ satisfying $\|w\|_2 =(\sum_{j=1}^{\infty}w_j^2)^{1/2}= O(j\log^3 n/\sqrt{n})=o(1)$, $w_j=0$ and $|w_k|\lesssim \Delta_\psi(j,k)$ for all $k\in\mathbb N\backslash\{j\}$ such that
\begin{equation}
A_n (e_j+w) = s_j(e_j+w),
\label{eq:eigen-an}
\end{equation}
and furthermore the convergence in the $o(1)$ term to zero as $n\to\infty$ is \emph{uniform} over $j\in[\bar J]$.
This shows that the self-adjoint operator $T_n=\sum_{k,\ell=1}^{n}[A_n]_{k\ell}\psi_k\otimes\psi_\ell$ admits a pair of eigenvalue and eigenfunction $(\hat s,\hat\psi)$ that satisfies
$\hat s=s_j(1+O(j\log^3 n/\sqrt{n}))=s_j(1+o(1))$, $\|\hat\psi-\psi_j\|_2=O(j\log^3 n/\sqrt{n})=o(1)$ and $|w_k|\lesssim \Delta_\psi(j,k)$ for all $k\neq j$. 
Furthermore:
\begin{enumerate}
\item These eigenfunctions of $T_n$ are unique for different $j$, because an $1-o(1)$ mass is concentrated on $e_j$; therefore, $\varpi$ is an injection;
\item Because $\hat s_j \geq (1+o(1))c_s j^{-2r}$ and $\hat s_j\leq (1+o(1))C_s j^{-2r}$, it holds that $j/C\leq \varpi(j)\leq Cj$;
\item All eigenvalue and eigenfunction pairs $(\hat s,\hat\psi)$ for $j\in[\bar J]$ will cover the top $J(n)$ eigenvalues and eigenfunctions of $T_n$ simply from the sufficiently large number of eigenvalues;
this shows that $\varpi$ is a surjection on $[J(n)]$.
\end{enumerate}

Let $H=A_n-A$ and consider $w_j=0$ which is simply a Gauge choice. Because $Ae_j=s_je_j$, Eq.~(\ref{eq:eigen-an}) can be reformulated as
$$
(A-s_j I+H)w = -He_j.
$$
Let $\bar H$ be the submatrix (also infinite-dimensional) formed by removing the $j$th row and $j$th column of $H$. Let $h$ be the $j$th column of $H$, removing $H_{jj}$.
Let $D=\diag(s_1-s_j,\cdots,s_{j-1}-s_j,s_{j+1}-s_j,\cdots)$ be a diagonal matrix of the same shape of $\bar H$. 
Let $\bar w=(w_1,\cdots,w_{j-1},w_{j+1},\cdots)$. Then
$$
\bar w = -(D+\bar H)^{-1}h = -D^{-1/2}(I+\tilde H)^{-1}D^{-1/2}h,
$$
where $\tilde H = D^{-1/2} \bar H D^{-1/2}$. 

For any $i,k\neq j$, Eq.~(\ref{eq:proof-stability-1}) implies
$$
\big|\tilde H_{ik}\big| \lesssim \frac{1}{\sqrt{n}} \frac{i^{-r}k^{-r}\log(nik)}{\sqrt{|(s_i-s_j)(s_k-s_j)|}} \leq \frac{1}{\sqrt{n}}z_iz_k\;\;\;\;\;\;\text{where}\;\; z_k = \frac{k^{-r}{\sqrt{\log n}\log(3k)}}{\sqrt{|s_k-s_j|}},
$$
and the $\lesssim$ inequality hides only constants and is uniform over all $i,k$.
In a similar way, $D^{-1/2}h$ can also be upper bounded component-wise as
\begin{equation}
\big|[D^{-1/2}h]_k \big| \lesssim  \frac{1}{\sqrt{n}}\frac{j^{-r}k^{-r}\log(njk)}{\sqrt{|s_k-s_j|}}\leq \frac{j^{-r}\log(nj)}{\sqrt{n}}z.
\label{eq:proof-stability-3}
\end{equation}

Discuss the following three cases:
\begin{enumerate}
\item The case of $k\leq j/C$. In this case, by Assumption (B1), $|s_k-s_j|\geq c_s k^{-2r} - C_s(Ck)^{-2r}\geq 0.5c_sk^{-2r}$, and therefore $|z_k|/\sqrt{\log n}\log(3k)\leq \sqrt{2/c_s}=O(1)$;
\item The case of $k\geq Cj$. In this case, by Assumption (B1), $|s_k-s_j| \geq c_s j^{-2r} - C_s (Cj)^{-2r}\geq 0.5c_s j^{-2r}$, and therefore $|z_k|/\sqrt{\log n}\log(3k)\leq \sqrt{2/c_s}j^r/k^r= O(j^rk^{-r})$.
\item The case of $j/C\leq k\leq Cj$. In this case, by Assumption (D1), $|s_k-s_j| \geq |k-j| \cdot c_g(Cj)^{-2r-1}$, and therefore 
$$
|z_k|/\sqrt{\log n}\log(3k)\leq \sqrt{C}j^{r+1/2}k^{-r}/\sqrt{c_g|k-j|}\leq C^{r+1}\sqrt{c_gk/|j-k|}= O(\sqrt{k/|j-k|}).
$$
\end{enumerate}
Because $\tilde H \leq zz^\top/\sqrt{n}$, summarizing the above cases we obtain 
\begin{align}
\|\tilde H\|_\op&\lesssim \frac{1}{\sqrt{n}}\|z\|_2^2 =\frac{1}{\sqrt{n}}\sum_{k\neq j}|z_k|^2 \nonumber\\
& \lesssim \frac{1}{\sqrt{n}}\left(j\log n\log^2(3j) + j^{2r}\int_j^{\infty}\frac{\log n\log^2 (3t)\ud t}{t^{2r}} + Cj\log n\log^2 (3Cj+3)\int_1^{Cj} \frac{\ud t}{t}\right)\nonumber\\
& \lesssim \frac{j\ln n\ln^2 (j+1))}{\sqrt{n}}\lesssim \frac{j\log^3 n}{\sqrt{n}}.
\end{align}
With $\bar J\asymp n^q$ for $q\in(0,1/2)$ (Because $J(n)\asymp n^q$ and $\bar J=CJ$, where $C$ is a constant), the right-hand side of the above inequality is $o(1)$ and therefore we have $\|\tilde H\|_\op=o(1)$ as $n\to\infty$.
We can then do the following Taylor expansion using $(I+\tilde H)^{-1} = \sum_{\ell=0}^{\infty} (-1)^\ell \tilde H^\ell$, for every $k\neq j$:
\begin{align}
|w_k| &\leq \sum_{\ell=0}^{\infty} \big|e_k^\top D^{-1/2}\tilde H^\ell D^{-1/2}h\big| \overset{(a)}{\leq }\sum_{\ell=0}^{\infty} \frac{j^{-r}\log(nj)}{\sqrt{n}}\left|e_k^\top D^{-1/2}\left(\frac{1}{\sqrt{n}}zz^\top\right)^\ell z\right|\nonumber\\
&\lesssim \sum_{\ell=0}^{\infty}\frac{j^{-r}\log(nj)}{\sqrt{n}}\left(\frac{j\log^3 n}{\sqrt{n}}\right)^\ell \big|e_k^\top D^{-1/2}z\big| \lesssim\frac{j^{-r}\log(nj)}{\sqrt{n}}\big|e_k^\top D^{-1/2}z\big|\nonumber\\
& = \frac{j^{-r}\log (nj)}{\sqrt{|s_k-s_j| n}}\big|z_k\big|.\label{eq:proof-stability-4} 
\end{align}
Here, Eq.~(a) holds by plugging in the expression of $D^{-1/2}h$ in Eq.~(\ref{eq:proof-stability-3}).
Invoking the case discussions earlier. For $k\leq j/C$, the right-hand side of Eq.~(\ref{eq:proof-stability-4}) is upper bounded by
$$
 \frac{j^{-r}\log(nj)}{\sqrt{|s_k-s_j|n}}\times \sqrt{\log n}\log(3k)\times O(1) = O\left(\frac{j^{-r}k^r\log^2 n\log(3j)\log(3k)}{\sqrt{n}}\right);
$$
for $k\geq Cj$, the right-hand side of Eq.~(\ref{eq:proof-stability-4}) is upper bounded by 
$$
\frac{j^{-r}\log(nj)}{\sqrt{|s_k-s_j|n}} \times\sqrt{\log n}\log(3k)\times O(j^rk^{-r}) = O\left(\frac{j^rk^{-r}\log^2 n\log(3j)\log(3k)}{\sqrt{n}}\right);
$$
for $j/C\leq k\leq Cj$, the right-hand side of Eq.~(\ref{eq:proof-stability-4}) is upper bounded by
\begin{align*}
\frac{j^{-r}\log(nj)}{\sqrt{|s_k-s_j|n}}\times O\left(\sqrt{\frac{k\log n\log^2(3k)}{|j-k|}}\right) &= \frac{j^{-r}\log^2 n}{\sqrt{n}}\times O\left(\frac{j^{r+1/2}}{\sqrt{|k-j|}}\right)\times O\left(\sqrt{\frac{k\log n\log^2(3j)}{|j-k|}}\right)\\
&= O\left(\frac{\sqrt{jk}\log^2 n\log(3j)\log(3k)}{|j-k|\sqrt{n}}\right).
\end{align*}
Combining all of the above cases we complete the proof of upper bounds on $|\langle\hat\psi_j,\psi_k\rangle|$ in Lemma \ref{lem:stability}.

Furthermore, 
\begin{align}
\|w\|_2 &\leq \frac{j^{-r}\log (nj)\sqrt{\log n}}{\sqrt{n}}\left(\int_1^j \frac{\log^3(3t)\ud t}{t^{-2r}} + j^{2r}\int_j^{\bar J}\frac{t^{-2r}\log^2(3t)}{j^{-2r}}\ud t \right.\nonumber\\
&\;\;\;\;\;\left.+ Cj\cdot \int_1^{Cj}\frac{1}{tj^{-2r-1}}\frac{\log^2(3t)\ud t}{t}\right)^{1/2}\nonumber\\
&\lesssim \frac{j^{-r}\log n\log^2(3Cj+3)}{\sqrt{n}}\left(j^{2r+1} + j^{2r+1} + j^{2r+2}\right)^{1/2} \lesssim  \frac{j\log^3 n}{\sqrt{n}}\leq \frac{\bar J\sqrt{\log n}}{\sqrt{n}}.\label{eq:proof-stability-5}
\end{align}
Because $\bar J\asymp n^q$ for some $q<1/2$, the right-hand side of Eq.~(\ref{eq:proof-stability-5}) is $o(1)$ uniformly for all $j\leq\bar J$. This completes the proof.

Finally, with the additional assumption that $q<1/4$, the right-hand side of Eq.~(\ref{eq:proof-stability-5}) is further upper bounded by $o(1/\bar J)=o(1/J(n))$. 
Additionally, in this case we have for every $j\leq \bar J$ that $\hat s_j-\hat s_{j+1} = s_j-s_{j+1}+o(s_j/\bar J)=s_j-s_{j+1}+o(j^{-2r}/\bar J) = (1+o(1))(s_j-s_{j+1})$ because $s_j-s_{j+1}\gtrsim j^{-2r-1}$.
This shows that there is a one-to-one match between $(s_j,\psi_j)$ and $(\hat s_j,\hat\psi_j)$ for all $j\leq\bar J$.



\subsection{Proof of Lemma \ref{lem:zpsi-stability}}
Fix arbitrary $z\in\supp(P_Z^n)$, $\ell\leq J$ and $\ell/C_J\leq j\leq C_J\ell$ such that $\varpi(j)=\ell$. 
Recall that for $a,b>0$, $(a+b)^M = \sum_{m=0}^M \binom{M}{m}a^mb^{M-m}$ from the binomial expansion property.
Therefore, for $M\in\{2,4\}$, 
\begin{align}
&\big|\langle z,\hat\psi_\ell\rangle^M-\langle z,\psi_j\rangle^M\big|\nonumber\\
& \leq \langle z,\psi_j\rangle^M(1-\langle\hat\psi_\ell,\psi_j\rangle^M) + 6\sum_{m=0}^{M-1}\big|\langle z,\psi_j\rangle\big|^m\left(\sum_{k\neq j}\big|\langle z,\psi_k\rangle\langle\hat\psi_\ell,\psi_k\rangle\big|\right)^{M-m}\nonumber\\
&\leq 2\langle z,\psi_j\rangle^M\sum_{k\neq j}\langle\hat\psi_\ell,\psi_k\rangle^2 +  6\sum_{m=0}^{M-1}\big|\langle z,\psi_j\rangle\big|^m\left(\sum_{k\neq j}\big|\langle z,\psi_k\rangle\langle\hat\psi_\ell,\psi_k\rangle\big|\right)^{M-m}\label{eq:zpsi-1}\\
&\leq 2L_J^M\left[\langle z,\psi_j\rangle^M\sum_{k\neq j}\Delta_\psi(j,k)^2 +3\sum_{m=0}^{M-1}\big|\langle z,\psi_j\rangle\big|^m\ \left(\sum_{k\neq j}\big|\langle z,\psi_k\rangle\big|\Delta_\psi(j,k)\right)^{M-m} \right].\label{eq:zpsi-2}
\end{align}
Here, Eq.~(\ref{eq:zpsi-1}) holds because $1-\langle\hat\psi_\ell,\psi_j\rangle^2=\sum_{k\neq j}\langle\hat\psi_\ell,\psi_k\rangle^2$ from the Pythagorean theorem,
and $1-\langle\hat\psi_\ell,\psi_j\rangle^4 = (1+\langle\hat\psi_\ell,\psi_j\rangle^2)(1-\langle\hat\psi_\ell,\psi_j\rangle^2) \leq 2(1-\langle\hat\psi_\ell,\psi_j\rangle^2)$. 
By Assumption (B3), $|\langle z,\psi_k\rangle|\leq C_\nu\sqrt{s_k}n^{\kappa} \lesssim k^{-r}n^{\kappa}$ almost surely.
Subsequently, the right-hand side of Eq.~(\ref{eq:zpsi-2}) is upper bounded almost surely, up to constants depending only on $q$ and other parameters in assumptions, by
\begin{align}
 n^{M\kappa}\left[j^{-rM}\sum_{k\neq j}\Delta_\psi(j,k)^2 + \sum_{m=0}^{M-1} j^{-rm}\left(\sum_{k\neq j}k^{-r}\Delta_\psi(j,k)\right)^{M-m} \right].
 \label{eq:zpsi-3}
\end{align}
Using expressions of $\Delta_\psi(\cdot,\cdot)$, we have that
\begin{align}
\sum_{k\neq j}\Delta_\psi(j,k)^2 &\lesssim \frac{\log^2 n\log^2(3j)}{n}\times \left( \frac{1}{j^{2r}}\int_1^{C_J j} k^{2r}\log^2(3k)\ud k \right.\nonumber\\
&\;\;\;\;\left.+ j^{2r}\int_{C_J j}^{\infty}\frac{\log^2(3k)\ud k}{k^{2r}} + 3C_J j^2\log^2(3C_Jj)\cdot \int_1^{C_j j}\frac{\ud t}{t^2}\right)\nonumber\\
&\lesssim \frac{j^2\log^6 n}{n},\label{eq:zpsi-4}
\end{align}
and 
\begin{align}
\sum_{k\neq j}k^{-r}\Delta_\psi(j,k) &\lesssim  \frac{\log n\log(3j)}{\sqrt{n}}\times \left( \frac{1}{j^{r}}\int_1^{C_J j}k^{-r}\cdot k^{r}\log(3k)\ud k \right.\nonumber\\
&\;\;\;\;\left.+ j^{r}\int_{C_J j}^{\infty}k^{-r}\cdot\frac{\log(3k)\ud k}{k^{r}} + \frac{(3C_J)^{1/2} j\log(3C_Jj)}{(j/C_j)^{r}}\cdot \int_1^{C_j j}\frac{\ud t}{t}\right)\nonumber\\
&\lesssim \frac{j^{1-r}\log^{4}n}{\sqrt{n}}.\label{eq:zpsi-5}
\end{align}
%\begin{align}
%\sum_{k\neq j}k^{-r\tau}\Delta_\psi(j,k)^\tau &\lesssim  \frac{\log^\tau n\log^\tau(3j)}{n^{\tau/2}}\times \left( \frac{1}{j^{r\tau}}\int_1^{C_J j}k^{-r\tau}\cdot k^{r\tau}\log^\tau(3k)\ud k \right.\nonumber\\
%&\;\;\;\;\left.+ j^{r\tau}\int_{C_J j}^{\infty}k^{-r\tau}\cdot\frac{\log^\tau(3k)\ud k}{k^{r\tau}} + \frac{(3C_J)^{\tau/2} j^\tau\log^\tau(3C_Jj)}{(j/C_j)^{r\tau}}\cdot \int_1^{C_j j}\frac{\ud t}{t^\tau}\right)\nonumber\\
%&\lesssim \frac{j^{\tau(1-r)}\log^{3\tau+1}n}{n^{\tau/2}}.\label{eq:zpsi-5}
%\end{align}
Subsequently,
\begin{align}
\sum_{m=0}^{M-1}& j^{-rm}\left(\sum_{k\neq j}k^{-r}\Delta_\psi(j,k)\right)^{M-m}
\lesssim \sum_{m=0}^{M-1}j^{-rm}\left(\frac{j^{1-r}\log^4 n}{\sqrt{n}}\right)^{M-m}\nonumber\\
&= \sum_{m=0}^{M-1}j^{-rM}\left(\frac{j\log^4 n}{\sqrt{n}}\right)^{M-m} \lesssim \frac{J^{1-rM}\log^4 n}{\sqrt{n}},
\label{eq:zpsi-6}
\end{align}
where the last inequality holds because $j\log^4n/\sqrt{n}\leq J\log^4 n/\sqrt{n} = O(n^{1/2-q}\log^4 n)\to 0$ and as $n\to\infty$, and therefore the $m=M-1$ term dominates.
Combining Eqs.~(\ref{eq:zpsi-3},\ref{eq:zpsi-4},\ref{eq:zpsi-6}) and noting that $j\lesssim n^q$ for some $q<1/2$, we complete the proof of the first statement in Lemma \ref{lem:zpsi-stability}.

Finally, because the right-hand side of Eq.~(\ref{eq:zpsi-2}) is a low-degree polynomial of $\{\langle z,\psi_k\rangle\}_{k=1}^{\infty}$,
to prove the second inequality involving expectations, simply replace Assumption (B3) with Assumption (B2) and remove the $n^{M\kappa}$ factor. 
%\pd{This last part is unclear. Suggest the following note.}
More specifically, {for the second inequality, set $z = Z$ in Eq.~\eqref{eq:zpsi-1} and take expectations. By Assumption (B2), $\mathbb{E}\left| \langle Z, \psi_k \rangle \right|^p \lesssim s^{p/2}_k$ for $p=2,4$. Combining these moment bounds with Cauchy Schwarz over $\sum_{k \neq j} \left| \langle Z, \psi_k \rangle \right| | \langle \hat\psi_\ell, \psi_k \rangle |$ and the bound established in Eq.~\eqref{eq:zpsi-3} we obtain the second inequality involving expectations, which is the same bound as the first inequality with the factor $n^{M\kappa}$ removed.}


\subsection{Technical lemmas for results in Secs.~\ref{sec:anitr} and \ref{sec:adptr}.}


\begin{lemma}
%Let $\{\psi_j\}_{j=1}^{\infty}$ and $\{\hat\psi_j\}_{j=1}^n$ be the eigenfunctions of $T$ and $T_n$, listed in descending order with respect to their eigenvalues.
Let $J=J(n)$ and $\und J=J(n)/C_J$.
Conditioned on results in Lemma \ref{lem:stability}, it holds that
%$$
%\|\hat\Pi_{J}^\perp z\|_2^2 \leq \|\Pi_{\und J}^\perp z\|_2^2 + L_J\times \sum_{j=1}^{\und J}\sum_{k\neq j}\langle z,\psi_k\rangle^2\Delta_\psi(j,k)^2\;\;\;\;\;\;\forall z\in L_2(I),
%$$
%. Furthermore, with $X\sim P_X^n$ and $Z=L_{K^{1/2}}X$, it holds that
$$
\mathbb E_{Z\sim P_Z^n}[\|\hat\Pi_J^\perp Z\|_2^2] \lesssim J^{1-2r}\left(1+\frac{J\log^4 n}{\sqrt{n}}\right),
$$
where $\hat\Pi_{J}^\perp = I-\sum_{j=1}^{J}\hat\psi_j\otimes\hat\psi_j$.
%and
%$$
%\|\hat\Pi_J^\perp Z\|_2^2 \;\;\lesssim \;\;n^{2\kappa}\times  J^{1-2r}\left(1+\frac{J\log^6 n}{n}\right) \;\;\;\;\;\;a.s.
%$$
\label{lem:truncation}
\end{lemma}

\begin{proof}[Proof of Lemma \ref{lem:truncation}.]
Fix arbitrary $z\in L_2(I)$ and, without loss of generality, assume $\|z\|_2=1$.
Throughout this proof, all expectations are taken with respect to $Z$.
The Pythagorean theorem asserts that
\begin{align}
\|z\|_2^2 = \|\hat\Pi_{J}z\|_2^2+\|\hat\Pi_{J}^\perp z\|_2^2 = \|\Pi_{\und J}z\|_2^2+\|\Pi_{\und J}^\perp z\|_2^2,
\label{eq:proof-truncation-1}
\end{align}
where $\hat\Pi_J=\sum_{j=1}^J\hat\psi_j\otimes\hat\psi_j$, $\Pi_{\und J}=\sum_{j=1}^{\und J}\psi_j\otimes\psi_j$ and $\hat\Pi_J^\perp=I-\hat\Pi_J$, $\Pi_{\und J}^\perp = I-\Pi_{\und J}$.
Using Lemmas \ref{lem:stability} and \ref{lem:zpsi-stability}, it holds that
\begin{align}
\mathbb E[\|\hat\Pi_J Z\|_2^2] &= \sum_{j=1}^J \mathbb E[\langle Z,\hat\psi_j\rangle^2] \geq \sum_{j=1}^{\und J}\mathbb E[\langle Z,\hat\psi_{\varpi(j)}\rangle^2]\nonumber\\
&\geq \left(\sum_{j=1}^{\und J}\mathbb E[\langle Z,\psi_j\rangle^2]\right) - O\left(\frac{J^{2-2r}\log^4 n}{\sqrt{n}}\right)\nonumber\\
&= \mathbb E[\|\Pi_{\und J} Z\|_2^2] -O\left(\frac{J^{2-2r}\log^4 n}{\sqrt{n}}\right) \label{eq:proof-truncation-2}
\end{align}
%\pd{
%\begin{align}
%\mathbb E[\|\hat\Pi_J^\perp Z\|_2^2] &= \mathbb E[\| Z\|_2^2]-\sum_{j=1}^J \mathbb E[\langle Z,\hat\psi_j\rangle^2] \nonumber\\
%&\leq  E[\| Z\|_2^2]-\sum_{j=1}^{\und J}\mathbb E[\langle Z,\psi_{j}\rangle^2]+\left| \sum_{j=1}^{\und J}\mathbb E[\langle Z,\psi_{j}\rangle^2] - \sum_{j=1}^J \mathbb E[\langle Z,\hat\psi_j\rangle^2]\right| \\ & \dots .\label{eq:proof-truncation-2}
%\end{align}
%}
Combining Eqs.~(\ref{eq:proof-truncation-1},\ref{eq:proof-truncation-2}) and noting that $\mathbb E[\|\Pi_{\und J}^\perp Z\|_2^2 \lesssim \int_{\und J}^{\infty} j^{-2r}\ud j \lesssim \bar J^{1-2r}$, 
we complete the proof of Lemma \ref{lem:truncation}, in the following manner:
{\color{magenta}
\begin{align*}
\mathbb E[\|\hat\Pi_J^\perp Z\|_2^2] = \mathbb E[\|\Pi_{\und J} Z\|_2^2] + \mathbb E[\Pi_{\und J}^\perp Z\|_2^2] - \mathbb E[\|\hat\Pi_J Z\|_2^2],
\end{align*}
where $\mathbb E[\|\Pi_{\und J} Z\|_2^2]$ cancels out the same term in the expression of Eq.~(\ref{eq:proof-truncation-2}).
}
%Invoking Lemma \ref{lem:stability}, for every $j\leq \und J$, there exists a unique $\ell\leq J$ such that $|\langle\hat\psi_\ell,\psi_j\rangle| = 1-o(1)\leq 1$, and
%$$
%\big|\langle\hat\psi_\ell, \psi_k\rangle\big| \leq L_J\Delta_\psi(j,k),\;\;\;\;\;\;\forall k\in\mathbb N,\;k\neq j.
%$$
%Subsequently,
%\begin{align}
%&\;\;\|\hat\Pi_J z\|_2^2 = \sum_{j=1}^J\langle z,\hat\psi_j\rangle^2 \geq \sum_{j=1}^{\und J}\langle z,\hat\psi_{\varpi(j)}\rangle^2 = \sum_{j=1}^{\und J}\left\langle z, \sum_{k=1}^{\infty}\langle\hat\psi_{\varpi(j)},\psi_k\rangle\psi_k\right\rangle^2\nonumber\\
%&\geq \sum_{j=1}^{\und J}\left[\langle z,\psi_j\rangle^2\langle\psi_j,\hat\psi_{\varpi(j)}\rangle^2 -  2\big|\langle z,\psi_j\rangle\big| \sum_{k\neq j}\big|\langle z,\psi_k\rangle\langle\psi_k,\hat\psi_{\varpi(j)}\rangle)\big|\right]\nonumber\\
%&\geq \sum_{j=1}^{\und J}\langle z,\psi_j\rangle^2 - \sum_{j=1}^{\und J}\left[\langle z,\psi_j\rangle^2\sum_{k\neq j}L_J^2\Delta_\psi(j,k)^2 + 2L_J\big|\langle z,\psi_j\rangle\big| \sum_{k\neq j}\big|\langle z,\psi_k\rangle\Delta_\psi(j,k)\big|\right]\label{eq:proof-truncation-1-1}\\
%&= \|\Pi_{\und J} z\|_2^2 - \sum_{j=1}^{\und J}\left[\langle z,\psi_j\rangle^2\sum_{k\neq j}L_J^2\Delta_\psi(j,k)^2 + 2L_J\big|\langle z,\psi_j\rangle\big| \sum_{k\neq j}\big|\langle z,\psi_k\rangle\Delta_\psi(j,k)\big|\right].\label{eq:proof-truncation-2}
%\end{align}
%Here, Eq.~(\ref{eq:proof-truncation-1-1}) holds because $\langle\psi_j,\hat\psi_{\varpi(j)}\rangle^2 = 1-\sum_{k\neq j}\langle\psi_k,\hat\psi_{\varpi(j)}\rangle^2 \geq 1-\sum_{k\neq j}L_J^2\Delta_\psi(j,k)^2$
%thanks to the Pythagorean theorem and Lemma \ref{lem:stability}.
%It is worth pointing out that Eq.~(\ref{eq:proof-truncation-2}) holds for any $z\in \supp(P_Z^n)$.
%
%Note that $\mathbb E[|\langle Z,\psi_j\rangle\langle Z,\psi_k\rangle|] \leq \sqrt{s_js_k}\lesssim j^{-r}k^{-r}$ for any $j,k\geq 1$, thanks to Assumption (B1). Subsequently,
%\begin{align}
%\sum_{j=1}^{\und J}&\mathbb E[\langle Z,\psi_j\rangle^2]\sum_{k\neq j}\Delta_\psi(j,k)^2\nonumber\\ %\leq \sum_{j>\und J} s_j+ L_J^2\times \sum_{j=1}^{\und J}\sum_{k\neq j}s_k\Delta_\Psi(j,k)^2\label{eq:proof-truncation-3}\\
%&\lesssim \sum_{j=1}^{\und J}\frac{1}{j^{2r}}\times \frac{\log^2 n\log^2(3j)}{n}\times \left( \frac{1}{j^{2r}}\int_1^{C_J j} k^{2r}\log^2(3k)\ud k \right.\nonumber\\
%&\;\;\;\;\left.+ j^{2r}\int_{C_J j}^{\infty}\frac{\log^2(3k)\ud k}{k^{2r}} + 3C_J j^2\log^2(3C_Jj)\cdot \int_1^{C_j j}\frac{\ud t}{t^2}\right) \nonumber\\
%&\lesssim \sum_{j=1}^{\und J}j^{2-2r}\times\frac{\log^6 n}{n}\lesssim \frac{J^{3-2r}\log^6 n}{n};\label{eq:proof-truncation-4}
%\end{align}
%additionally,
%\begin{align}
%\sum_{j=1}^{\und J}& \mathbb E\left[\big|\langle Z,\psi_j\rangle\big|\sum_{k\neq j}\big|\langle Z,\psi_k\rangle|\Delta_\psi(j,k)\right]\nonumber\\
%&\lesssim \sum_{j=1}^{\und J}\frac{1}{j^r}\times \frac{\log n\log(3j)}{\sqrt{n}}\times \left( \frac{1}{j^{r}}\int_1^{C_J j} k^{-r}\cdot k^{r}\log(3k)\ud k \right.\nonumber\\
%&\;\;\;\;\left.+ j^{r}\int_{C_J j}^{\infty}k^{-r}\cdot\frac{\log(3k)\ud k}{k^{r}} + \frac{\sqrt{3C_J }j\log(3C_Jj)}{(j/C_j)^r}\cdot \int_1^{C_j j}\frac{\ud t}{t}\right) \nonumber\\
%&\lesssim \sum_{j=1}^{\und J}j^{1-2r}\times \frac{\log^4 n}{\sqrt{n}}\lesssim \frac{J^{2-2r}\log^4 n}{\sqrt{n}}.\label{eq:proof-truncation-5}
%\end{align}
%Combining Eqs.~(\ref{eq:proof-truncation-1},\ref{eq:proof-truncation-2},\ref{eq:proof-truncation-5}) and noting that $\mathbb E[\|\Psi_{\und J}^\perp Z\|_2^2 \lesssim \int_{\und J}^{\infty} j^{-2r}\ud j \lesssim \bar J^{1-2r}$, we complete the proof of Lemma \ref{lem:truncation}.
%
%To prove the second inequality, note that $\mathbb E[\langle Z,\psi_k\rangle^4]\leq C_p s_k^2$ for every $k$, thanks to Assumption (B2).
%
%
%Finally, for the last inequality, note that both Eqs.~(\ref{eq:proof-truncation-3},\ref{eq:proof-truncation-4}) would hold almost surely for $Z$, if both right-hand sides are multiplied by $C_\nu^2 n^{2\kappa}$.
%\pd{We need to state somewhere that all expectations are with respect to $Z$ independent of the sample.}
\end{proof}


The following lemma holds for \emph{any} choices of $\Lambda_J$.
\begin{lemma}
Let $J=J(n)$ and suppose $q\in(0,1/2)$ and $\kappa<(1-q)/2$.
Conditioned on the results in Lemma \ref{lem:truncation}, it holds that 
$$
\sup_{x\in \supp(P_X^n)}\frac{\sum_{i=1}^n |g(z,Z_i)|^3}{(\sum_{i=1}^n g(z,Z_i)^2)^{3/2}} \;\;\lesssim \;\; \frac{\sqrt{J}n^\kappa}{\sqrt{n}}\left(1+\frac{\sqrt{J}\log^2 n}{\sqrt[4]{n}}\right) \;\; \to 0\;\;\;\;\;\;\text{as }n\to\infty,
$$
where $z=L_{K^{1/2}}x$, $Z_i=L_{K^{1/2}}X_i$ and $g(z,Z)=\langle z, (T_J+\Lambda_J)^\dagger Z\rangle$.
\label{lem:anitr-lyapnov}
\end{lemma}
\begin{proof}[Proof of Lemma \ref{lem:anitr-lyapnov}.]
Let $\varsigma^2 = \frac{1}{n}\sum_{i=1}^n g(z,Z_i)^2 = \langle z, (T_J+\Lambda_J)^\dagger T_n (T_J+\Lambda_J)^\dagger z\rangle = \langle z,(T_J+\Lambda_J)^{\dagger}T_J(T_J+\Lambda_J)^\dagger z\rangle 
= \|T_J^{1/2}(T_J+\Lambda_J)^\dagger z\|_2^2$.
Note that here in the second equality, the middle operator $T_n$ is switched to $T_J$, which is valid because $(T_J+\Lambda_J)^\dagger$
is fully within the span of $\{\hat\psi_j\}_{j=1}^J$ by the definition of Moore-penrose pseudo inverses on operators.
% \pd{In the second equality, $T_n$ has been switch to $T_J$ but the error needs to be accounted for. We can just lower bound $\langle z, (T_J+\Lambda_J)^\dagger T_n (T_J+\Lambda_J)^\dagger z\rangle$ with a constant multiple of $\|T_J^{1/2}(T_J+\Lambda_J)^\dagger z\|_2^2$ and then move forward.}
Using H\"{o}lder's inequality, we have that
\begin{align}
\frac{\sum_{i=1}^n |g(z,Z_i)|^3}{(\sum_{i=1}^n g(z,Z_i)^2)^{3/2}} &\leq \frac{\max_{1\leq i\leq n}|g(z,Z_i)|}{\sqrt{n}\varsigma}
= \max_{1\leq i\leq n}\frac{|\langle z, (T_J+\Lambda_J)^\dagger Z_i\rangle|}{\sqrt{n}\|T_J^{1/2}(T_J+\Lambda_J)^\dagger z\|_2}\nonumber\\
&\leq \max_{1\leq i\leq n}\frac{\|T_J^{\dagger/2} Z_i\|_2}{\sqrt{n}},\label{eq:proof-anitr-lyapnov-1}
\end{align}
where Eq.~(\ref{eq:proof-anitr-lyapnov-1}) holds by using the Cauchy-Schwarz inequality on $\langle z, (T_J+\Lambda_J)^\dagger Z_i\rangle$ to cancel the demoninator,
and $T_J^{\dagger/2} = (T_J^\dagger)^{1/2}=\sum_{j=1}^J (\hat s_j)^{-1/2}\hat\psi_j\otimes\hat\psi_j$.

By the results of Lemma \ref{lem:truncation}, for any $\ell\leq J$, there exists $j\leq \bar J=C_J J$ such that $\varpi(j)=\ell$. We denote $\varpi^{-1}(\ell)=j$ in this case. 
Using also Lemma \ref{lem:zpsi-stability}, it holds almost surely that
\begin{align}
\|T_J^{\dagger/2}Z_i\|_2^2  &= \sum_{\ell=1}^J\frac{\langle Z_i,\hat\psi_\ell\rangle^2}{\hat s_\ell}
\leq \sum_{\ell=1}^J\frac{2\langle Z_i,\hat\psi_\ell\rangle^2}{s_{\varpi^{-1}(\ell)}}\nonumber\\
&\leq \sum_{\ell=1}^J \frac{2}{s_{\varpi^{-1}(\ell)}}\left[\langle Z_i,\psi_{\varpi^{-1}(\ell)}\rangle^2 + O\left(\frac{n^{2\kappa}J^{1-2r}\log^4 n}{\sqrt{n}}\right)\right]\nonumber\\
&\lesssim \sum_{j=1}^{\bar J}O(n^{2\kappa})\times\left[1 + \frac{1}{s_j}\cdot \frac{J^{1-2r}\log^4 n}{\sqrt{n}}\right]\nonumber\\
&\lesssim Jn^{2\kappa} + J^2 n^{2\kappa}\log^4 n/\sqrt{n},\label{eq:proof-anitr-lyapnov-2}
\end{align}
where the last inequality holds because $\sum_{j=1}^{\bar J} {\frac{1}{s_j}} \lesssim \int_1^{\bar J}j^{2r}\ud j \lesssim \bar J^{1+2r}\lesssim J^{1+2r}$.
Combining Eqs.~(\ref{eq:proof-anitr-lyapnov-1},\ref{eq:proof-anitr-lyapnov-2}) we complete the proof of Lemma \ref{lem:anitr-lyapnov}.
%\begin{align}
%\|T_J^{\dagger/2}Z_i\|_2^2 &= \sum_{\ell=1}^J\frac{\langle Z_i,\hat\psi_\ell\rangle^2}{\hat s_\ell}
%= \sum_{\ell=1}^J\sum_{k=1}^{\infty}\frac{\langle Z_i,\psi_k\rangle^2\langle\hat\psi_\ell,\psi_k\rangle^2}{\hat s_\ell}\nonumber\\
%&\leq \sum_{\ell=1}^J \left[\frac{\langle Z_i,\psi_{\varpi^{-1}(\ell)}\rangle^2}{\hat s_\ell} + \sum_{k\neq\varpi^{-1}(\ell)}\frac{\langle Z_i,\psi_k\rangle^2\langle\hat\psi_\ell,\psi_k\rangle^2}{\hat s_\ell} \right]\nonumber\\
%&\leq \sum_{\ell=1}^J \left[\frac{\langle Z_i,\psi_{\varpi^{-1}(\ell)}\rangle^2}{\hat s_\ell} + L_J^2\sum_{k\neq\varpi^{-1}(\ell)}\frac{\langle Z_i,\psi_k\rangle^2\Delta_\psi(k,\varpi^{-1}(\ell))^2}{\hat s_\ell} \right]\label{eq:proof-anitr-lyapnov-2}\\
%&\leq C_\nu^2n^{2\kappa}\times \sum_{\ell=1}^J \left[\frac{s_{\varpi^{-1}(\ell)}}{\hat s_\ell} + L_J^2\sum_{k\neq\varpi^{-1}(\ell)}\frac{s_k\Delta_\psi(k,\varpi^{-1}(\ell))^2}{\hat s_\ell} \right]\label{eq:proof-anitr-lyapnov-3}\\
%&\leq 2C_\nu^2n^{2\kappa}\times \sum_{\ell=1}^J \left[1 + \frac{L_J^2}{s_{\varpi^{-1}(\ell)}}\sum_{k\neq\varpi^{-1}(\ell)}s_k\Delta_\psi(k,\varpi^{-1}(\ell))^2 \right].\label{eq:proof-anitr-lyapnov-4}
%\end{align}
%Here, Eq.~(\ref{eq:proof-anitr-lyapnov-2}) holds by applying Lemma \ref{lem:stability}; 
%Eq.~(\ref{eq:proof-anitr-lyapnov-3}) holds almost surely thanks to Assumption (B3);
%Eq.~(\ref{eq:proof-anitr-lyapnov-4}) holds by applying Lemma \ref{lem:stability} again on the upper bound of $|\hat s_\ell/s_{\varpi^{-1}(\ell)}-1|$.
%For any $j$ and $\ell=\varpi(j)$, Assumption (B1) and the fact that $j/C_J\leq \ell\leq C_J j$ yield
%\begin{align}
%\frac{1}{s_j}\sum_{k\neq j}s_k\Delta_\psi(j,k)^2 &\lesssim j^{2r}\times \frac{\log^2 n\log^2(3j)}{n}\times \left(\int_1^{j/C_J}k^{-2r}\cdot\frac{k^{2r}\log^2(3k)\ud k}{j^{2r}}\right. \nonumber\\
%&\;\;\;\;\left.+ j^{2r}\int_{C_J j}^{\infty}k^{-2r}\cdot \frac{\log^2(3k)\ud k}{k^{2r}} + \frac{(C_J j)^2\log^2(3C_J j)}{(j/C_j)^{2r}}\int_1^{C_J j}\frac{\ud t}{t^2}\right)^2\nonumber\\
%&\lesssim \frac{j^2\log^6 n}{n}.\nonumber
%\end{align}
%Subsequently, 
%\begin{equation}
%\|T_J^{\dagger/2}Z_i\|_2^2 \lesssim Jn^{2\kappa}\left(1 + \frac{J^2\log^6 n}{n}\right).
%\label{eq:proof-anitr-lyapnov-5}
%\end{equation}
%Combining Eqs.~(\ref{eq:proof-anitr-lyapnov-1},\ref{eq:proof-anitr-lyapnov-5}) and we complete the proof of Lemma \ref{lem:anitr-lyapnov}.
\end{proof}

\begin{lemma}
There exists a constant $n_0$ depending only on parameters on assumptions, such that for all $n\geq n_0$,
conditioned on results in Lemma \ref{lem:zpsi-stability}, the following inequalities hold:
\begin{align*}
\mathbb E_{P_Z^n}[\sigma(Z)^2] &\leq 2\sum_{j=1}^{J(n)}\frac{s_{\varpi^{-1}(j)}}{(s_{\varpi^{-1}(j)}+\lambda_{nj})^2}\left[s_{\varpi^{-1}(j)} + O\left(\frac{J(n)^{1-2r}\log^4 n}{\sqrt{n}}\right)\right];\\
\mathbb E_{P_Z^n}[\sigma(Z)^2] &\geq \frac{1}{2}\sum_{j=1}^{J(n)}\frac{s_{\varpi^{-1}(j)}}{(s_{\varpi^{-1}(j)}+\lambda_{nj})^2}\left[s_{\varpi^{-1}(j)} + O\left(\frac{J(n)^{1-2r}\log^4 n}{\sqrt{n}}\right)\right];\\
\mathbb E_{P_Z^n}[\sigma(Z)^4] &\leq 2\sum_{j=1}^{J(n)}\frac{s_{\varpi^{-1}(j)}^2}{(s_{\varpi^{-1}(j)}+\lambda_{nj})^4}\left[C_p s_{\varpi^{-1}(j)}^2 + O\left(\frac{J(n)^{1-4r}\log^4 n}{\sqrt{n}}\right)\right]
\end{align*}
where $\varpi^{-1}$ is its inverse map of $\varpi$ implied in Lemma \ref{lem:truncation}, which exists on $[J(n)]$ on which $\varpi$ is a surjection,
and $\sigma(Z)^2 = \langle Z, (T_J+\Lambda_J)^\dagger T_J(T_J+\Lambda_J)^\dagger Z\rangle$.
Constants in all $O(\cdot)$ notations depend only on parameters in assumptions.
\label{lem:variance-stability}
\end{lemma}

\begin{remark}
If Assumptions (B1) and (C2) hold with $q<1/2$, then all terms in big-O notations in Lemma \ref{lem:variance-stability} can be safely dropped, at a cost of perhaps larger $n_0$ that still depends only on parameters in assumptions.
\end{remark}

\begin{proof}[Proof of Lemma \ref{lem:variance-stability}.]
Let $J=J(n)$.
It holds that
\begin{align*}
\sigma(Z)^2 &= \sum_{j=1}^J \frac{\langle Z,\hat\psi_j\rangle^2\hat s_j}{(\hat s_j+\lambda_{nj})^2} = (1+o(1))\sum_{j=1}^J\frac{\langle Z,\hat\psi_j\rangle^2 s_{\varpi^{-1}(j)}}{(s_{\varpi^{-1}(j)}+\lambda_{nj})^2};\\
\sigma(Z)^4 &= \sum_{j=1}^J\frac{\langle Z,\hat\psi_j\rangle^4\hat s_j^2}{(\hat s_j+\lambda_{nj})^4} = (1+o(1))\sum_{j=1}^J\frac{\langle Z,\hat\psi_j\rangle^4 s_{\varpi^{-1}(j)}^2}{(s_{\varpi^{-1}(j)}+\lambda_{nj})^4}.
\end{align*} 
\pd{Fourth moment are missing the cross terms in the expansion. Those need to be accounted for. Since we need to only upper bound this later, needs a Cauchy Schwarz and moment-control bound to repair this.}

{\color{red} YW: this is a known issue. The entire upper bound on the truncation term (high probablity bound) won't work. Let's discuss, maybe in conjunction with the additional assumption needed for the previous Lyapnov condition issue.}

All inequalities in Lemma \ref{lem:variance-stability} are then consequences of Lemma \ref{lem:zpsi-stability}.
\end{proof}

\begin{lemma}
Impose all assumptions in Lemmas \ref{lem:stability} and \ref{lem:zpsi-stability}.
Suppose for some $a>0$, $\lambda_n=o(n^{-a})$ and $\lambda_n=\omega(n^{-a-\delta})$ for any $\delta>0$.
If $\frac{a}{2r}<q<\frac{1}{2}$, then for any $\epsilon>0$, there exist $n_\epsilon\in\mathbb N$ and $C_\epsilon<+\infty$, both depending only on $\epsilon,a$ and other parameters in assumptions, such that
$$
\Pr_{Z\sim P_Z^n}\left[\sum_{j=1}^{J(n)}\frac{\langle Z,\hat\psi_j\rangle^2\hat s_j}{(\hat s_j+2\sqrt{\hat s_j\lambda_n})^2} \geq C_\epsilon \lambda_n^{-1/2r}\right]\;\geq\;1-\epsilon
$$
for every $n\geq n_\epsilon$ and $P_Z^n$, and the statement is conditioned on the success event in Lemma \ref{lem:zpsi-stability} over $\{\hat s_j,\hat\psi_j\}$.
\label{lem:truncation-whp}
\end{lemma}
\begin{proof}[Proof of Lemma \ref{lem:truncation-whp}.]
Let $J=J(n)$ and $\bar J=C_JJ(n)$.
Let $\sigma(Z)^2 = \langle z, (T_J+\Lambda_J)^\dagger T_J(T_J+\Lambda_J)^\dagger z\rangle \pd{\langle Z, (T_J+\Lambda_J)^\dagger T_J(T_J+\Lambda_J)^\dagger Z\rangle}$, where $T_J=\sum_{j=1}^J\hat s_j\hat\psi_j\otimes\hat\psi_j$
and $\Lambda_J = \sum_{j=1}^J 2\sqrt{\hat s_j\lambda_n}\hat\psi_j\otimes\hat\psi_j$.
Using Chebyshev's inequality, for any $\epsilon>0$ it holds with probability $1-\epsilon/2$ that
\begin{align}
\sigma(Z)^2 &\geq \mathbb E[\sigma(Z)^2] - \sqrt{\frac{2}{\epsilon}}\sqrt{\mathbb E[\sigma(Z)^4]}.
\label{eq:proof-anitr-limiting-7}
\end{align}
For any $j\leq J$, Lemma \ref{lem:stability} implies that, for every $t,\tau\geq 1$, 
$$
\frac{s_{\varpi^{-1}(j)}^t}{(s_{\varpi^{-1}(j)}+2(1+o(1))\sqrt{s_{\varpi^{-1}(j)}\lambda_n})^\tau} \asymp \frac{j^{-2rt}}{j^{-2r\tau}+j^{-r\tau}\lambda_n^{\tau/2}},
$$
where the $\asymp$ notation hide only constants depending on assumption parameters, and hold uniformly for all $j$ and $n$. Lemma \ref{lem:variance-stability} then yields,
with the condition $q<1/2$, that
\begin{align*}
\mathbb E[\sigma(Z)^2] &\gtrsim \int_1^J\frac{j^{-4r}\ud j}{j^{-4r}+j^{-2r}\lambda_n}\overset{(a)}{\gtrsim} \lambda_n^{-1/2r}  \gtrsim  \lambda_n^{-1/2r};\\
\mathbb E[\sigma(Z)^4] &\lesssim \int_1^J\frac{j^{-8r}\ud j}{j^{-8r}+j^{-4r}\lambda_n^2}
\lesssim \lambda_n^{-1/2r} \lesssim \lambda_n^{-1/2r}.
\end{align*}
Here, Eq.~(a) holds if $q>a/2r$, because $\lambda_n =o(n^{-a})$ and $\lambda_n=\omega(n^{-a-\delta})$ for any $\delta>0$.
With $a>0$, this shows that $\mathbb E[\sigma(Z)^2]\gtrsim\lambda_n^{-1/2r}$ and $\sqrt{\mathbb E[\sigma(Z)^4]}\lesssim \lambda_n^{-1/4r} \ll \lambda_n^{-1/2r}$. \pd{This part needs to be fixed.}
%With $q<\frac{1}{4}+\frac{a}{4r}$, the above inequalities imply that $\mathbb E[\sigma(Z)^2]\gtrsim \lambda_n^{-1/2r}$ and $\sqrt{\mathbb E[\sigma(Z)^4]}\lesssim \lambda_n^{-1/4r} \ll \lambda_n^{-1/2r}$.
Hence, Eq.~(\ref{eq:proof-anitr-limiting-7}) implies that, for some constant $C'<+\infty$, with probability $1-\epsilon/2$ 
\begin{equation*}
\sigma(Z)^2 \geq C'\lambda_n^{-1/2r},
\end{equation*}
which is to be proved.
%\begin{align*}
%\mathbb E[\sigma(Z)^2] &\gtrsim \int_1^J\frac{j^{-4r}\ud j}{j^{-4r}+j^{-2r}\lambda_n} - O\left(\frac{J^{1-2r}\log^4 n}{\sqrt{n}}\right)\times \int_1^J\frac{j^{-2r}\ud j}{j^{-4r}+j^{-2r}\lambda_n}\nonumber\\
%&\overset{(a)}{\gtrsim} \lambda_n^{-1/2r} - O\left(\frac{J^{1-2r}\log^4 n}{\sqrt{n}}\right)\times J^{1+2r} \gtrsim  \lambda_n^{-1/2r}- O\left(\frac{J^2\log^4 n}{\sqrt{n}}\right);\\
%\mathbb E[\sigma(Z)^4] &\lesssim \int_1^J\frac{j^{-8r}\ud j}{j^{-8r}+j^{-4r}\lambda_n^2} + O\left(\frac{J^{1-4r}\log^4 n}{\sqrt{n}}\right)\times \int_1^J\frac{j^{-4r}\ud j}{j^{-8r}+j^{-4r}\lambda_n^2}\nonumber\\
%&\lesssim \lambda_n^{-1/2r}+ O\left(\frac{J^{1-4r}\log^4 n}{\sqrt{n}}\right)\times J^{1+4r} \lesssim \lambda_n^{-1/2r} + O\left(\frac{J^2\log^4 n}{\sqrt{n}}\right).
%\end{align*}
%Here, Eq.~(a) holds if $q>a/2r$, because $\lambda_n =o(n^{-a})$ and $\lambda_n=\omega(n^{-a-\delta})$ for any $\delta>0$.
%With $q<\frac{1}{4}+\frac{a}{4r}$, the above inequalities imply that $\mathbb E[\sigma(Z)^2]\gtrsim \lambda_n^{-1/2r}$ and $\sqrt{\mathbb E[\sigma(Z)^4]}\lesssim \lambda_n^{-1/4r} \ll \lambda_n^{-1/2r}$.
%Hence, Eq.~(\ref{eq:proof-anitr-limiting-7}) implies that, for some constant $C'<+\infty$, with probability $1-\epsilon/2$ 
%\begin{equation*}
%\sigma(Z)^2 \geq C'\lambda_n^{-1/2r},
%\end{equation*}
%which is to be proved.
\end{proof}

\subsection{Proof of Theorem \ref{thm:anitr-limiting}}

Let $J=J(n)$, $T_J=\sum_{j=1}^J \hat s_j\hat\psi_j\otimes \hat\psi_j$ and $\Lambda_J=\sum_{j=1}^J\lambda_{nj}\hat\psi_j\otimes\hat\psi_j$.
Decompose the estimation error of $\hat\theta_n^\anitr-\theta_0$ as
\begin{equation}
\hat\theta_n^\anitr-\theta_0 = \underbrace{-\hat\Pi_J^\perp\theta_0}_{=:t_n} \underbrace{-(T_J+\Lambda_J)^\dagger \Lambda_J\theta_0}_{=: b_n} + \underbrace{\frac{1}{n}\sum_{i=1}^n\varepsilon_i (T_J+\Lambda_J)^\dagger Z_i}_{=:v_n},
\label{eq:anitr-decomp}
\end{equation}
where $\hat\Pi_J^\perp = I - \sum_{j=1}^J\hat\psi_j\otimes\hat\psi_j$.

The rest of the proof is carried out fixing an arbitrary $x\in\supp(P_X)$ and $z=L_{K^{1/2}}x$. 
The truncation error is precisely $t_n(z)=\langle z, t_n\rangle = -\langle z, \hat\Pi_J^\perp\theta_0\rangle = -\langle \hat\Pi_J^\perp \pd{z},\theta_0\rangle$, 
because $\hat\Pi_J^\perp$ is a self-disjoint, projection operator.

For the bias term $b_n(z)=\langle z,b_n\rangle$, define $\hat\nu_j := \langle z,\hat\psi_j\rangle$ and $\hat\vartheta_j := \langle \theta_0,\hat\psi_j\rangle$ for all $j\leq J$. 
Recall also the definition that $\lambda_{nj} = \sqrt{\hat s_j\lambda_n}$ Then
\begin{align}
\big|b_n(z)\big|^2 &= \left|\sum_{j=1}^J\frac{\lambda_{nj}\hat\vartheta_j\hat\nu_j}{\hat s_j+\lambda_{nj}}\right|^2 \leq \sum_{j=1}^J\frac{\lambda_{nj}^2\hat\nu_j^2}{(\hat s_j+\lambda_{nj})^2}
=\lambda_n\sum_{j=1}^J\frac{\hat\nu_j^2\hat s_j}{(\hat s_j+\sqrt{\hat s_j\lambda_n})^2}\label{eq:proof-anitr-limiting-2}\\
&= \lambda_n\langle z, (T_J+\Lambda_J)^\dagger T_J(T_J+\Lambda_J)^\dagger z\rangle.\label{eq:proof-anitr-limiting-3}
\end{align}
Here, the second inequality in Eq.~(\ref{eq:proof-anitr-limiting-2}) holds by using the Cauchy-Schwarz inequality and noting that $\sum_{j=1}^J\hat\vartheta_j^2\leq \|\theta_0\|_2^2\leq 1$.
Because $\lambda_n n\to0$ as $n\to\infty$, Eq.~(\ref{eq:proof-anitr-limiting-3}) implies
\begin{equation}
\sup_{x\in \supp(P_X^n)}b_n(z)/\hat\sigma_n^\anitr(z)\;\;\to\;\;0\;\;\;\;\;\;a.s.
\label{eq:proof-anitr-limiting-3half}
\end{equation}
as $n\to\infty$, where $z=L_{K^{1/2}}x$.

For the variance term $v_n$, let $v_n(z)= \langle v_n,z\rangle = \frac{1}{n}\sum_{i=1}^n\varepsilon_i \langle z, (T_J+\Lambda_J)^\dagger Z_i\rangle$.
Using Lemma \ref{lem:anitr-lyapnov} and the finite-sample Berry-Esseen theorem for independently but not identically distributed random variables, there exists a constant $C<+\infty$
depending only on parameters of Assumptions, such that
\begin{equation}
\sup_{x\in\supp(P_X^n)} \|\hat\sigma_n^\anitr(z)^{-1}v_n(z) - u\|_{\ks} \;\;\leq\;\; C\left(n^{-c} + \frac{\sqrt{J}n^{\kappa}}{\sqrt{n}} \left(1+\frac{\sqrt{J}\log^2 n}{\sqrt[4]{n}}\right)\right),
\label{eq:anitr-limiting-4}
\end{equation}
where $z=L_{K^{1/2}}x$, $u\sim\mathcal N(0,1)$ and $\|P-Q\|_\ks := \sup_{t\in\mathbb R}|F_P(t)-F_Q(t)|$ is the Kolmogorov-Smirnov (KS) distance between two uni-variate distributions $P$ and $Q$,
with $F_P$ and $F_Q$ being their CDFs respectively.
The first term in Eq.~(\ref{eq:anitr-limiting-4}) accounts for the failure probability in Lemma \ref{lem:stability}.
Under the conditions imposed, the right-hand side of Eq.~(\ref{eq:anitr-limiting-4}) converges to zero as $n\to\infty$.
Combining Eqs.~(\ref{eq:proof-anitr-limiting-3half},\ref{eq:anitr-limiting-4}) we complete the proof of the first statement in Theorem \ref{thm:anitr-limiting}.

To prove the strengthened second statement in Theorem \ref{thm:anitr-limiting} under the additional conditions, first note that, using the Markov inequality, the second property of Lemma \ref{lem:truncation} implies
for every $\epsilon>0$ that 
\begin{equation}
\Pr_{Z\sim P_Z^n}\left[\|\hat\Pi_J^\perp Z\|_2 > \frac{CJ^{\frac{1}{2}-r}}{\epsilon}\left(1+\frac{\sqrt{J}\log^2 n}{\sqrt[4]{n}}\right)\right] \;\leq\; \epsilon/2,
\label{eq:proof-anitr-limiting-6}
\end{equation}
where $C<+\infty$ is a constant depending only on constants in the assumptions.
On the other hand, note that 
$$
\hat\sigma_n^\anitr(z)^2 = \frac{\sigma_\epsilon^2}{n}\langle z,(T_J+\Lambda_J)^\dagger T_J(T_J+\Lambda_J)^\dagger z\rangle \geq \frac{\sigma_\epsilon^2}{n}  \sum_{j=1}^J\frac{\langle z,\hat\psi_j\rangle^2\hat s_j}{(\hat s_j+2\sqrt{\hat s_j\lambda_n})^2}
$$
where the second inequality holds because $\lambda_{nj}=\lambda_n\wedge\sqrt{\hat s_j\lambda_n}\leq 2\sqrt{\hat s_j\lambda_n}$.
Invoking Lemma \ref{lem:truncation-whp} with $a=1$ and $\frac{1}{2r}<q<\frac{1}{2}$, for any $\epsilon>0$ it holds for sufficiently large $n$ that
\begin{equation}
\Pr_{Z\sim P_Z^n}\left[\hat\sigma_n^\anitr(Z) > \frac{C_\epsilon'\sigma_\varepsilon}{\sqrt{n}}\lambda_n^{-1/4r}\right]\geq 1-\epsilon/2,
\label{eq:proof-anitr-limiting-6half}
\end{equation}
where $C_\epsilon'<+\infty$ only depends on $\epsilon$ and other assumption parameters.
Comparing Eqs.~(\ref{eq:proof-anitr-limiting-6},\ref{eq:proof-anitr-limiting-6half}) and using the union bound, with probability $1-\epsilon$ it holds that
$$
\frac{\|\hat\Pi_J^\perp Z\|}{\hat\sigma_n^\anitr(Z)} \lesssim \sqrt{n}\lambda_n^{1/4r} J^{\frac{1}{2}-r} \to 0, 
$$
provided that $q>1/2r$. This completes the proof.


\subsection{Proof of Theorem \ref{thm:anitr-rates}}
Let $J=J(n)\asymp n^q$ for some $q\in(0,1/2)$.
Without explicitly mentioned, all expectations and probabilities in this proof are with respect to the randomness in $Z\sim P_Z^n$.
Invoking Lemma \ref{lem:truncation}, it holds that $\mathbb E[\|\hat\Pi_J^\perp Z\|_2^2]\lesssim J^{1-2r} \lesssim n^{-q(2r-1)}$. This explains the second term in Theorem \ref{thm:anitr-rates}.

Now let us focusing only on the $\sigma(Z)^2=\langle Z,(T_J+\Lambda_J)^\dagger T_J (T_J+\Lambda_J)^\dagger Z\rangle$ term in the definition of $\hat\sigma_n^\anitr(Z)$.
Using spectral representations, it holds that
\begin{equation}
\mathbb E[\sigma(Z)^2] = \sum_{j=1}^J \frac{\hat s_j \mathbb E[\langle Z,\hat\psi_j\rangle^2]}{(\hat s_j+\lambda_{nj})^2}
=\sum_{j=1}^J\frac{\mathbb E[\langle Z,\hat\psi_j\rangle^2]}{{\hat s_j}+\lambda_n}.
\label{eq:proof-anitr-rates-1}
\end{equation}
Applying Lemmas \ref{lem:stability} and \ref{lem:zpsi-stability}, with probability $1-O(n^{-c})$ it holds for sufficiently large $n$ that
\begin{align}
\sum_{j=1}^J&\frac{\mathbb E[\langle Z,\hat\psi_j\rangle^2]}{{\hat s_j}+{\lambda_n}}
\leq \sum_{j=1}^J\frac{2}{{s_{\varpi^{-1}(j)}}+\lambda_n}\left[s_{\varpi^{-1}(j)} + O\left(\frac{J^{1-2r}\log^4 n}{\sqrt{n}}\right)\right]\nonumber\\
&\lesssim \int_1^{\infty} \frac{j^{-2r}\ud j}{j^{-2r}+{\lambda_n}} + O\left(\frac{J^{1-2r}\log^4 n}{\sqrt{n}}\right)\times \int_1^{C_J j}\frac{\ud j}{j^{-2r}}\nonumber\\
&\lesssim \lambda_n^{-1/2r} + O\left(\frac{J^2\log^4 n}{\sqrt{n}}\right).
\end{align}
With $J\asymp n^q$ and $\lambda_n=\omega(n^{-1-\delta})$ for any $\delta>0$, we complete the proof of Theorem \ref{thm:anitr-rates}.


\subsection{Proof of Theorem \ref{thm:adptr-limiting}}

The analysis of the truncation and variance terms are identical to the proof in Theorem \ref{thm:anitr-limiting}.
In the rest of this proof, we shall therefore focus on the bias term primarily, and additional constraints leading to the strengthened statement dropping truncation terms.

Let $J=J(n)$ and $\bar J=C_JJ$.
Fix arbitrary $x\in\supp(P_X^n)$ and let $z=L_{K^{1/2}}x$. Let $b_n(z)=\langle z, b_n\rangle$ be the bias term that is represented as
$$
b_n(z) = -\sum_{j=1}^J\frac{\lambda_{nj}\hat\vartheta_j\hat\nu_j}{\hat s_j+\lambda_{nj}},
$$
where $\hat\vartheta_j=\langle\theta_0,\hat\psi_j\rangle$ and $\hat\nu_j=\langle z, \hat\psi_j\rangle$. 
%It should be noted that the choice of
%$\lambda_{nj}=\hat s_j\sqrt{\mu_n}/|\hat\nu_j|$ if $\hat\nu_j^2\geq 2\hat s_j$ and $\lambda_{nj}=\sqrt{\hat s_j\mu_n}$ otherwise also depends on $z$
% \pd{This is not consistent with Eq. (17). For $\hat\nu_j^2< 2\hat s_j$, it should be $\lambda_{nj}=\mu_n \wedge \sqrt{\hat s_j \mu_n}$. Is it sufficint to take $\lambda_{nj}= \sqrt{\hat s_j \mu_n}$ for $\hat\nu_j^2< 2\hat s_j$?}
Using H\"{o}lder's inequality and Assumption (D1), it holds almost surely that
\begin{align}
\big|b_n(z)\big| \leq \left(\sum_{j=1}^{J}|\vartheta_j|\right)\times \max_{1\leq j\leq J}\frac{\lambda_{nj}|\hat\nu_j|}{\hat s_j+\lambda_{nj}}.
\label{eq:proof-adptr-limiting-1}
\end{align}
Recall the definition that $\theta_k=\langle\theta_0,\psi_k\rangle$ for every $k\geq 1$.
%Because $q<1/4$, the map $\varpi$ in Lemma \ref{lem:stability} can be chosen as the bijection $\varpi(j)=j$. 
{\color{magenta}
Because $q<1/2$, the map $\varpi$ in Lemma \ref{lem:stability} can be chosen such that it is a surjection on $[J]$, meaning that for every $j\in[J]$, there exists a unique $\ell\leq C_J j$
such that $\varpi(\ell)=j$. For notational simplicity let $\varrho$ be the inverse function of $\varpi$ on $[J]$, so that $\varpi(\varrho(j))=j$ for all $j\leq J$.
%\pd{Is the bijection property needed or can it be waived?}
Conditioned on the results in Lemma \ref{lem:stability}, it holds for every $j\leq J$ that
\begin{align}
\big|\hat\vartheta_j\big| &= \big|\langle \theta_0,\hat\psi_{\varrho(j)}\rangle\big| \leq \big|\langle \theta_0,\psi_{\varrho(j)}\rangle\big| + \sum_{k\neq \varrho(j)}\big|\langle\theta_0,\psi_k\rangle\langle\hat\theta_\ell,\psi_k\rangle\big|
=\big|\theta_j\big| + \sum_{k\neq \varrho(j)}\big|\theta_k\langle\hat\psi_\ell,\psi_k\rangle\big|\nonumber\\
&\leq \big|\theta_{\varrho(j)}\big| + \sqrt{\sum_{k\neq \varrho(j)}\langle\hat\psi_\ell,\psi_k\rangle^2} \leq \big|\theta_{\varrho(j)}\big| + L_J\sqrt{\sum_{k\neq {\varrho(j)}}\Delta_\psi(\varrho(j),k)^2}\label{eq:proof-adptr-limiting-2}\\
&\leq |\theta_{\varrho(j)}| + O\left(\frac{j\log^3 n}{\sqrt{n}}\right)\label{eq:proof-adptr-limiting-3}
\end{align}
where the first inequality in Eq.~(\ref{eq:proof-adptr-limiting-2}) holds by using the Cauchy-Schwarz inequality and noting that $\theta_0\in\Theta_1(\mH)\subseteq\Theta_2(\mH)$ implies $\sum_{j=1}^{\infty}\theta_j^2<+\infty$,
and the second inequality in Eq.~(\ref{eq:proof-adptr-limiting-2}) holds by applying Lemma \ref{lem:stability}; the last inequality holds by plugging in the expressions of $\Delta_\psi(\varrho(j),k)$,
noting that $\varrho(j)\leq C_J j=O(j)$
and following the analysis that leads to Eq.~(\ref{eq:zpsi-4}) in the proof of Lemma \ref{lem:zpsi-stability}.
With $J(n)\asymp n^q$ for $q\in(0,1/2)$, both $\varrho$ and $\varpi$ are injections and subsequently Eqs.~(\ref{eq:proof-adptr-limiting-1},\ref{eq:proof-adptr-limiting-3}) imply that
\begin{equation}
\big|b_n(z)\big| \leq C\times  \max_{1\leq j\leq J}\frac{\lambda_{nj}|\hat\nu_j|}{\hat s_j+\lambda_{nj}}
\label{eq:proof-adptr-limiting-4}
\end{equation}
for some $C>0$ depending only on assumption parameters. 
}

For every $j\leq J$, conduct a case analysis: 
\begin{enumerate}
\item If $\hat\nu_j^2\geq 2\hat s_j$, then $\lambda_{nj}=\hat s_j\sqrt{\mu_n}/|\hat\nu_j|$ and therefore
$$
\frac{\lambda_{nj}|\hat\nu_j|}{\hat s_j+\lambda_{nj}}\leq \frac{\hat s_j\sqrt{\mu_n}}{\hat s_j+\lambda_{nj}}\leq \sqrt{\mu_n};
$$
\item If $\hat\nu_j^2< 2\hat s_j$, then $\lambda_{nj}=\sqrt{\hat s_j\mu_n}$ and therefore
$$
\frac{\lambda_{nj}|\hat\nu_j|}{\hat s_j+\lambda_{nj}} = \frac{\sqrt{\hat s_j\mu_n}\sqrt{2\hat s_j}}{\hat s_j+\sqrt{\hat s_j\mu_n}} = \frac{\sqrt{2\mu_n}}{1+\sqrt{\mu_n/\hat s_j}}\leq \sqrt{2\mu_n};
$$
\end{enumerate}
Combining the above inequalities with Eq.~(\ref{eq:proof-adptr-limiting-4}), we have
$$
\big|b_n(z)\big| \leq C\times  \sqrt{\mu_n},
$$
which proves the upper bound on $|b_n(z)|$.

Next, we show that under additional conditions $\omega_n$ can be set to zero. Note that, with the choice of $\lambda_{nj}(z)$, it always holds that $\lambda_{nj}(z)\leq 2\sqrt{\hat s_j\mu_n}$.
Subsequently,
$$
\hat\sigma_n^\adptr(z)^2 \geq \frac{\sigma_\varepsilon^2}{n}\sum_{j=1}^J \frac{\langle z,\hat\psi_j\rangle^2\hat s_j}{(\hat s_j+2\sqrt{\hat s_j\mu_n})^2}.
$$
Invoking Lemma \ref{lem:truncation-whp} with $a=2r/(2r+1)$, under the conditions that $\frac{1}{2r+1}<q<\frac{1}{2}$, for any $\epsilon>0$ it holds with probability at least $1-\epsilon/2$ for all sufficiently large $n$ that
$$
\hat\sigma_n^\adptr(z)^2 \gtrsim n^{-1}\mu_n^{-1/2r}.
$$
Combining this with Eq.~(\ref{eq:proof-anitr-limiting-6}) in the proof of Theorem \ref{thm:anitr-limiting}, with probability $1-\epsilon$ it holds for sufficiently large $n$ that
$$
\frac{\|\hat\Pi_J^\perp Z\|_2+|b_n(z)|}{\hat\sigma_n^\adptr(Z)} \lesssim \sqrt{n}\mu_n^{1/4r}J^{1/2-r} + \sqrt{n}\mu_n^{1/4r}\sqrt{\mu_n} \to 0,
$$
provided that $q>2r/((2r-1)(2r+1))$. 

%\pd{If bijection is not needed here, then we might get a more generous range of $r$ ad then it works for $r$ such that $2r/((2r-1)(2r+1)) < 1/2$, that is, $r > \frac{1+\sqrt{2}}{2}$.}



\subsection{Proof of Theorem \ref{thm:adptr-rates}}

The $n^{-q(2r-1)+\delta}$ and $n^{1-c}$ terms arise from the truncation errors and failure probabilities of Lemma \ref{lem:stability}, which enjoy the same analysis and upper bounds
in the proof of Theorem \ref{thm:anitr-rates}.
Furthermore, $\omega_n\mu_n\lesssim n^{-\frac{2r}{2r+1}+\delta}$ for any $\delta>0$ simply from the asymptotic scalings of $\omega_n$ and $\mu_n$.
Therefore, we shall focus entirely on the $\hat\sigma_n^\adptr(z)^2$ term in the rest of this proof.

Let $J=J(n)$.
It holds that
\begin{align}
\hat\sigma_n^\adptr(z)^2 &\leq \frac{\sigma_\varepsilon^2}{n}\sum_{j=1}^J\frac{\hat\nu_j(z)^2\hat s_j}{(\hat s_j+\lambda_{nj}(z))^2} \leq \frac{\sigma_\varepsilon^2}{n}\sum_{j=1}^J\frac{\hat\nu_j(z)^2}{\hat s_j+\lambda_{nj}(z)}\nonumber\\
&\leq \frac{\sigma_\varepsilon^2}{n}\sum_{j=1}^J\frac{\hat\nu_j(z)^2}{\hat s_j+\sqrt{\hat s_j\mu_n}} + \frac{\sigma_\varepsilon^2}{n}\sum_{j=1}^J\frac{\hat\nu_j(z)^2}{\hat s_j+\hat s_j\sqrt{\mu_n}/|\hat\nu_j(z)|} \label{eq:proof-adptr-rates-0}\\
&\leq  \frac{\sigma_\varepsilon^2}{n}\underbrace{\sum_{j=1}^J\frac{\hat\nu_j(z)^2}{\hat s_j+\sqrt{\hat s_j\mu_n}}}_{=:\mA(z)} + \frac{\sigma_\varepsilon^2}{n}\underbrace{\sum_{j=1}^J\frac{|\hat\nu_j(z)|^3}{\hat s_j(|\hat\nu_j(z)|+\sqrt{\mu_n})}}_{=:\mB(z)}.
\label{eq:proof-adptr-rates-1}
\end{align}
where Eq.~(\ref{eq:proof-adptr-rates-0}) by splitting the two case of $\hat\nu_j^2(z) \geq 2\hat s_j$, $\hat\nu_j^2(z) <2\hat s_j$ and including both of them in the upper bound.

For term $\mA$, conditioned on the results of Lemma \ref{lem:stability} with $\varrho$ being the inverse function of injection $\varpi$ on $[J]$, we upper bound its expectation as
{\color{magenta}
\begin{align}
\mathbb E_{Z\sim P_Z^n}[\mA(Z)] &\leq \sum_{j=1}^J \frac{\mathbb E[\hat\nu_j(Z)^2]}{\hat s_j + \mu_n} \leq \sum_{j=1}^J\frac{2\mathbb E[\hat\nu_j(Z)^2]}{s_{\varrho(j)}+\mu_n}.\label{eq:proof-adptr-sp1}
\end{align}
for sufficiently large $n$.
}

For term $\mB$, we upper bound its expectation as
\begin{align}
\mathbb E_{Z\sim P_Z^n}[\mB(Z)] &\leq \sum_{j=1}^{J}\frac{\sqrt{\mathbb E\left[\hat\nu_j(Z)^4\right]}}{\hat s_j}\sqrt{\mathbb E\left[\frac{\hat\nu_j(Z)^2}{(|\hat\nu_j(Z)|+\sqrt{\mu_n})^2}\right]}\label{eq:proof-adptr-rates-2}\\
&\leq\sum_{j=1}^{J}\frac{\sqrt{\mathbb E\left[\hat\nu_j(Z)^4\right]}}{\hat s_j}\sqrt{\mathbb E\left[\frac{\hat\nu_j(Z)^2}{\hat\nu_j(Z)^2+\mu_n}\right]} \label{eq:proof-adptr-rates-3}\\
&\leq \sum_{j=1}^{J}\frac{\sqrt{\mathbb E\left[\hat\nu_j(Z)^4\right]}}{\hat s_j}\sqrt{\frac{\mathbb E[\hat\nu_j(Z)^2]}{\mathbb E[\hat\nu_j(Z)^2]+\mu_n}}\label{eq:proof-adptr-rates-3half}\\
&\leq \sum_{j=1}^{J}\frac{2\sqrt{\mathbb E\left[\hat\nu_j(Z)^4\right]}}{s_{\varrho(j)}}\sqrt{\frac{\mathbb E[\hat\nu_j(Z)^2]}{\mathbb E[\hat\nu_j(Z)^2]+\mu_n}}\label{eq:proof-adptr-rates-4}
\end{align}
Here, Eq.~(\ref{eq:proof-adptr-rates-2}) holds by Cauchy-Schwarz inequality;
Eq.~(\ref{eq:proof-adptr-rates-3}) holds by the fact that $(a+b)^2\geq a^2+b^2$ {for $a b \geq 0$};
Eq.~(\ref{eq:proof-adptr-rates-3half}) holds from the Jensen's inequality and the fact that $f(x)=x/(a+x)$ is a concave function of $x$ on $\mathbb R_+$ for every $a>0$;
Eq.~(\ref{eq:proof-adptr-rates-5}) holds by applying Lemma \ref{lem:stability} for sufficiently large $n$.

Next, conditioned on results in Lemma \ref{lem:zpsi-stability} and the fact that $\varpi(j)=j$ since $q<1/4$, for every $j\leq J$ we have
\begin{align}
\mathbb E[\hat\nu_j(Z)^2] &= \mathbb E[\langle Z,\psi_j\rangle^2] + O\left(\frac{J^{1-2r}\log^4 n}{\sqrt{n}}\right) \in [0.5s_j,2s_j];\label{eq:proof-adptr-rates-5}\\
\mathbb E[\hat\nu_j(Z)^4] &\leq \mathbb E[\langle Z,\psi_j\rangle^4] + O\left(\frac{J^{1-2r}\log^4 n}{\sqrt{n}}\right)\leq 2C_p s_j^2,\label{eq:proof-adptr-rates-5half}
\end{align}
where the second inequality in Eq.~(\ref{eq:proof-adptr-rates-5half}) holds thanks to Assumption (B2).
Because $s_J\gtrsim J^{-2r}$, with the condition that 

Combining Eqs.~(\ref{eq:proof-adptr-rates-1}), (\ref{eq:proof-adptr-sp1}), (\ref{eq:proof-adptr-rates-4}), (\ref{eq:proof-adptr-rates-5}), and (\ref{eq:proof-adptr-rates-5half}), we obtain
\begin{align*}
\mathbb E_{Z\sim P_Z^n}\left[\hat\sigma_n^\adptr(Z)]^2\right]
&\leq \frac{\sigma_\varepsilon^2}{n}\left(\sum_{j=1}^J\frac{4s_j}{s_{\varrho(j)}+\mu_n} + \sum_{j=1}^J\frac{2\sqrt{C_p}s_{\varrho(j)}^{3/2}}{0.5s_{\varrho(j)}(\sqrt{s_{\varrho(j)}}+\sqrt{\mu_n})}\right)\\
&\lesssim \frac{1}{n}\left(\int_1^{\infty}\frac{j^{-2r}\ud j}{j^{-2r}+\mu_n} + \int_1^{\infty}\frac{j^{-r}\ud j}{j^{-r}+\sqrt{\mu_n}}\right)\\
&\lesssim \frac{\mu_n^{-1/2r}}{n}\lesssim n^{-\frac{2r}{2r+1}+\delta},
\end{align*}
for any $\delta>0$, {\color{magenta}where the second inequality holds because $j/C_J\leq\varrho(j)\leq C_J j$ implying that $\varrho(j)$ is on the same order of $j$ up to a numerical constant of $C_J$
that is uniform for all $j$.}
This proves the theorem.


\bibliographystyle{apalike}
\bibliography{refs}

\end{document}